% =====================================================================
%  TEMA · MATEMÁTICAS — variante del design system (BEAMER 16:9)
%  Identidad para decks de matemáticas (azul \setsubject{mat}):
%    · FIX del bug de preflight §5/§6: las láminas a sangre del design
%      system (\portada \seccion \idea \pregunta \accion) usan
%      remember picture + current page y salen CORRIDAS hacia arriba con
%      compile_quiet.sh. Aquí se redefinen con `background canvas` +
%      coordenadas absolutas (misma técnica probada del tema eco/cnat).
%    · carga lmodern/colortbl/enumitem (que design/beamer.tex no trae) para
%      math escalable en titulares, \rowcolor en tablas y itemize[clave=valor].
%    · \DisplaySB con BoldFont explícito → \textbf dentro de \idea no avisa.
%    · portada con el personaje "Dr. Cuack" matemático (opcional; se dibuja
%      solo si existe el PNG).
%
%  INVOCACIÓN (tras % =====================================================================
%  DESIGN SYSTEM BEAMER v5 — ITSTZ · 16:9
%  Coherente con el libro: Inter Display (títulos) + Inter (cuerpo),
%  un color de asignatura, fondo claro, texto GRANDE y contrastado.
%  Filosofía slide: una idea por lámina, mucho aire, cero relleno.
%
%  Vocabulario:
%    \portada{kicker}{TÍTULO}{subtítulo}
%    \seccion{texto}                  lámina divisoria a pantalla completa
%    \idea{TEXTO GRANDE}              una afirmación que llena la slide
%    \resalta{texto}                  realce inline (fondo color)
%    frames de contenido normales con \frametitle
%    \begin{lineah} \hito{x}{etiqueta} ... \end{lineah}   timeline horizontal
%    \begin{lineav} \hitov{x}{título}{texto} ... \end{lineav} timeline vertical
%    \begin{mapa}{NODO CENTRAL} \rama{ang}{dist}{texto} ... \end{mapa}
%    \pregunta{texto}                 pregunta viva (apertura)
%    \accion{texto}                   cierre "de la queja a la acción"
%  Uso: % =====================================================================
%  DESIGN SYSTEM BEAMER v5 — ITSTZ · 16:9
%  Coherente con el libro: Inter Display (títulos) + Inter (cuerpo),
%  un color de asignatura, fondo claro, texto GRANDE y contrastado.
%  Filosofía slide: una idea por lámina, mucho aire, cero relleno.
%
%  Vocabulario:
%    \portada{kicker}{TÍTULO}{subtítulo}
%    \seccion{texto}                  lámina divisoria a pantalla completa
%    \idea{TEXTO GRANDE}              una afirmación que llena la slide
%    \resalta{texto}                  realce inline (fondo color)
%    frames de contenido normales con \frametitle
%    \begin{lineah} \hito{x}{etiqueta} ... \end{lineah}   timeline horizontal
%    \begin{lineav} \hitov{x}{título}{texto} ... \end{lineav} timeline vertical
%    \begin{mapa}{NODO CENTRAL} \rama{ang}{dist}{texto} ... \end{mapa}
%    \pregunta{texto}                 pregunta viva (apertura)
%    \accion{texto}                   cierre "de la queja a la acción"
%  Uso: % =====================================================================
%  DESIGN SYSTEM BEAMER v5 — ITSTZ · 16:9
%  Coherente con el libro: Inter Display (títulos) + Inter (cuerpo),
%  un color de asignatura, fondo claro, texto GRANDE y contrastado.
%  Filosofía slide: una idea por lámina, mucho aire, cero relleno.
%
%  Vocabulario:
%    \portada{kicker}{TÍTULO}{subtítulo}
%    \seccion{texto}                  lámina divisoria a pantalla completa
%    \idea{TEXTO GRANDE}              una afirmación que llena la slide
%    \resalta{texto}                  realce inline (fondo color)
%    frames de contenido normales con \frametitle
%    \begin{lineah} \hito{x}{etiqueta} ... \end{lineah}   timeline horizontal
%    \begin{lineav} \hitov{x}{título}{texto} ... \end{lineav} timeline vertical
%    \begin{mapa}{NODO CENTRAL} \rama{ang}{dist}{texto} ... \end{mapa}
%    \pregunta{texto}                 pregunta viva (apertura)
%    \accion{texto}                   cierre "de la queja a la acción"
%  Uso: \input{design/beamer.tex} ; \setsubject{csoc} ; \begin{document}...
% =====================================================================
\documentclass[aspectratio=169, 11pt]{beamer}
\usetheme{default}
\setbeamertemplate{navigation symbols}{}   % sin botones de navegación

\usepackage{fontspec}
\usepackage{polyglossia}\setdefaultlanguage{spanish}
\usepackage{tikz}
\usetikzlibrary{calc, positioning, arrows.meta, decorations.pathreplacing}
\usepackage{amsmath}

% --- TIPOGRAFÍA (idéntica familia que el libro) -----------------------
\setsansfont{Inter}[UprightFont=Inter-Regular, BoldFont=Inter-SemiBold,
  ItalicFont=Inter-Italic, Ligatures=TeX, Numbers=Lining]
\newfontfamily\Display{Inter Display}[Ligatures=TeX]
\newfontfamily\DisplayBlk{Inter Display Black}[Ligatures=TeX]
\newfontfamily\DisplaySB{Inter Display SemiBold}[Ligatures=TeX]
\newfontfamily\DisplayThin{Inter Display Thin}[Ligatures=TeX]

% --- COLOR (mismas 4 paletas que el libro) ----------------------------
\definecolor{mat} {RGB}{ 30,  64, 120}
\definecolor{cnat}{RGB}{ 16, 110,  80}
\definecolor{csoc}{RGB}{165,  75,  45}
\definecolor{eco} {RGB}{ 70,  55, 130}
\definecolor{ink}{RGB}{28,28,30}        % casi-negro para texto
\definecolor{paper}{RGB}{250,250,248}   % blanco cálido de fondo
\colorlet{subjectcolor}{csoc}
% \setsubject re-aplica los colores de beamer además de actualizar el alias
\newcommand{\setsubject}[1]{%
  \colorlet{subjectcolor}{#1}%
  \setbeamercolor{structure}{fg=#1}%
  \setbeamercolor{frametitle}{fg=#1,bg=}%
}

\setbeamercolor{normal text}{fg=ink, bg=paper}
\setbeamercolor{structure}{fg=subjectcolor}
\setbeamerfont{frametitle}{family=\DisplaySB}
\setbeamercolor{frametitle}{fg=subjectcolor, bg=}
\setbeamertemplate{frametitle}{%
  \vspace{0.6em}\hspace{0.2em}%
  {\DisplaySB\large\color{subjectcolor}\insertframetitle}\par
  \vspace{-0.2em}{\color{subjectcolor!30}\rule{\dimexpr\paperwidth-1.4em\relax}{1.2pt}}}

% Pie minimal: instituto · profe · tema · número en cuadro de color
\newcommand{\bmrtema}{}
\newcommand{\bmrprofe}{}   % ← \renewcommand{\bmrprofe}{Prof. Nombre} en cada beamer
\setbeamertemplate{footline}{%
  \hbox{}\vspace{2pt}\hspace{0.6em}%
  {\footnotesize\color{ink!45}ITSTZ%
   \ifx\bmrprofe\empty\else\space·\space\bmrprofe\fi}%
  \hfill{\footnotesize\color{ink!45}\bmrtema}\hspace{0.6em}%
  \tikz[baseline=-3pt]{\node[fill=subjectcolor, text=white, inner xsep=6pt,
     inner ysep=2.5pt, font=\DisplaySB\scriptsize]{\insertframenumber};}%
  \hspace{0.6em}\vspace{4pt}}

% =====================================================================
%  PORTADA y DIVISORES
% =====================================================================
\newcommand{\portada}[3]{% {kicker}{TÍTULO}{subtítulo}
{\setbeamertemplate{footline}{}
\begin{frame}[plain]
\begin{tikzpicture}[remember picture, overlay]
  \fill[paper] (current page.south west) rectangle (current page.north east);
  \node[anchor=north west, text=subjectcolor!12, font=\DisplayBlk]
    at ([xshift=-0.5cm,yshift=1.2cm]current page.north west)
    {\fontsize{190}{190}\selectfont ·};
  \node[anchor=west, text=subjectcolor, font=\DisplaySB\Large]
    at ([xshift=1.2cm, yshift=1.6cm]current page.west) {#1};
  \node[anchor=west, text=ink, font=\DisplayBlk, align=left]
    at ([xshift=1.15cm, yshift=0.2cm]current page.west)
    {\fontsize{40}{44}\selectfont #2};
  \fill[subjectcolor] ([xshift=1.2cm, yshift=-1.4cm]current page.west)
    rectangle ++(3.2cm, 0.16cm);
  \node[anchor=west, text=ink!60, font=\Display, align=left, text width=0.7\paperwidth]
    at ([xshift=1.2cm, yshift=-2.3cm]current page.west) {#3};
\end{tikzpicture}
\end{frame}}}

% Lámina divisoria: solo una frase grande centrada sobre color
\newcommand{\seccion}[1]{%
{\setbeamertemplate{footline}{}
\begin{frame}[plain]
\begin{tikzpicture}[remember picture, overlay]
  \fill[subjectcolor] (current page.south west) rectangle (current page.north east);
  \node[text=white, font=\DisplaySB, align=center, text width=0.85\paperwidth]
    at (current page.center) {\fontsize{30}{36}\selectfont #1};
\end{tikzpicture}
\end{frame}}}

% =====================================================================
%  ELEMENTOS DE CONTENIDO
% =====================================================================
% Una afirmación gigante que llena la lámina
\newcommand{\idea}[1]{%
{\setbeamertemplate{footline}{}
\begin{frame}[plain]
\begin{tikzpicture}[remember picture, overlay]
  \fill[paper] (current page.south west) rectangle (current page.north east);
  \fill[subjectcolor] (current page.north west)
    rectangle ([yshift=-0.32cm]current page.north east);
  \node[text=ink, font=\DisplaySB, align=center,
        text width=0.82\paperwidth]
    at (current page.center)
    {\fontsize{32}{44}\selectfont #1};
\end{tikzpicture}
\end{frame}}}

% Realce inline (fondo color suave, texto oscuro)
\newcommand{\resalta}[1]{\tikz[baseline=(R.base)]{\node[fill=subjectcolor!18,
  rounded corners=2pt, inner xsep=3pt, inner ysep=1pt](R){#1};}}
% Realce fuerte (fondo sólido, texto blanco)
\newcommand{\marca}[1]{\tikz[baseline=(R.base)]{\node[fill=subjectcolor,
  text=white, rounded corners=2pt, inner xsep=4pt, inner ysep=1.5pt](R){\textbf{#1}};}}

% Pregunta viva de apertura
\newcommand{\pregunta}[1]{%
\begin{frame}[plain]
\begin{tikzpicture}[remember picture,overlay]
  \fill[subjectcolor!7] (current page.south west) rectangle (current page.north east);
  \fill[subjectcolor] (current page.south west)
    rectangle ([xshift=0.45cm]current page.north west);
\end{tikzpicture}
\vfill
\hspace{1.5cm}\begin{minipage}{0.78\paperwidth}
  {\color{subjectcolor}\DisplayBlk\fontsize{54}{54}\selectfont ¿}\par
  \vspace{0.25em}
  {\Display\color{ink}\fontsize{22}{28}\selectfont #1}\par
  \vspace{0.4em}
  {\hfill\color{subjectcolor}\DisplayBlk\fontsize{38}{38}\selectfont ?}
\end{minipage}
\vfill
\end{frame}}

% Cierre / acción — label opcional (default: "DE LA QUEJA A LA ACCIÓN")
% Uso: \accion{texto}  o  \accion[ETIQUETA CUSTOM]{texto}
\newcommand{\accion}[2][DE LA QUEJA A LA ACCIÓN]{%
\begin{frame}[plain]
\begin{tikzpicture}[remember picture, overlay]
  \fill[subjectcolor!10] (current page.south west) rectangle (current page.north east);
  \node[anchor=west, text=subjectcolor, font=\DisplaySB]
    at ([xshift=1cm, yshift=2.2cm]current page.west)
    {\large #1};
  \node[anchor=west, text=ink, font=\Display, align=left, text width=0.78\paperwidth]
    at ([xshift=1cm, yshift=0cm]current page.west)
    {\fontsize{20}{26}\selectfont #2};
\end{tikzpicture}
\end{frame}}

% =====================================================================
%  LÍNEA DE TIEMPO HORIZONTAL
%  \begin{lineah}{año_ini}{año_fin}
%     \hito{posición 0-1}{etiqueta}   ...
%  \end{lineah}
% =====================================================================
\newenvironment{lineah}{%
  \begin{center}\begin{tikzpicture}[x=0.92\paperwidth, font=\footnotesize\sffamily]
  \draw[line width=2pt, subjectcolor, -{Stealth[length=5pt]}] (0,0) -- (1.02,0);
}{%
  \end{tikzpicture}\end{center}}
% hito: #1 = posición relativa 0..1 ; #2 = etiqueta (texto corto)
\newcommand{\hito}[2]{%
  \fill[subjectcolor] (#1,0) circle (3.2pt);
  \node[above, text=ink, align=center, text width=2.6cm, anchor=south,
        yshift=2pt] at (#1,0) {\bfseries #2};}
% hito bajo la línea (para alternar y no chocar)
\newcommand{\hitob}[2]{%
  \fill[subjectcolor] (#1,0) circle (3.2pt);
  \node[below, text=ink, align=center, text width=2.6cm, anchor=north,
        yshift=-2pt] at (#1,0) {#2};}

% =====================================================================
%  LÍNEA DE TIEMPO VERTICAL (ideal para listas cronológicas largas)
%  \begin{lineav} \hitov{año}{Título}{detalle} ... \end{lineav}
% =====================================================================
\newcounter{lvrow}
\newenvironment{lineav}{%
  \setcounter{lvrow}{0}%
  \begin{list}{}{\leftmargin=1.6cm \itemsep=0.55em \topsep=0.3em}%
  \renewcommand{\hitov}[3]{\stepcounter{lvrow}%
    \item[\tikz{\node[circle, fill=subjectcolor, text=white, inner sep=0pt,
       minimum size=1.5em, font=\DisplaySB\scriptsize]{\thelvrow};}]%
    {\DisplaySB\color{subjectcolor}##1}\quad{\bfseries ##2}\\[1pt]
    {\small\color{ink!75}##3}}%
}{%
  \end{list}}
\newcommand{\hitov}[3]{}  % placeholder; redefinido dentro del entorno

% =====================================================================
%  MAPA MENTAL — nodo central + ramas radiales
%  \begin{mapa}{NODO CENTRAL}
%     \rama{ángulo}{distancia cm}{texto de la rama}  ...
%  \end{mapa}
% =====================================================================
\newenvironment{mapa}[2][1]{%   % [escala opcional, def=1]{NODO CENTRAL}
  \begin{center}\begin{tikzpicture}[font=\sffamily\small, scale=#1, transform shape]
  \node[circle, fill=subjectcolor, text=white, font=\DisplaySB,
        align=center, text width=2.4cm, inner sep=4pt, minimum size=2.6cm]
        (CENTRO) at (0,0) {#2};
  \renewcommand{\rama}[3]{%
    \node[draw=subjectcolor!55, fill=subjectcolor!8, rounded corners=4pt,
          text=ink, align=center, text width=2.7cm, inner sep=4pt]
          (R##1) at (##1:##2) {##3};
    \draw[subjectcolor!55, line width=1.2pt] (CENTRO) -- (R##1);}%
}{%
  \end{tikzpicture}\end{center}}
\newcommand{\rama}[3]{}  % placeholder; redefinido dentro del entorno
 ; \setsubject{csoc} ; \begin{document}...
% =====================================================================
\documentclass[aspectratio=169, 11pt]{beamer}
\usetheme{default}
\setbeamertemplate{navigation symbols}{}   % sin botones de navegación

\usepackage{fontspec}
\usepackage{polyglossia}\setdefaultlanguage{spanish}
\usepackage{tikz}
\usetikzlibrary{calc, positioning, arrows.meta, decorations.pathreplacing}
\usepackage{amsmath}

% --- TIPOGRAFÍA (idéntica familia que el libro) -----------------------
\setsansfont{Inter}[UprightFont=Inter-Regular, BoldFont=Inter-SemiBold,
  ItalicFont=Inter-Italic, Ligatures=TeX, Numbers=Lining]
\newfontfamily\Display{Inter Display}[Ligatures=TeX]
\newfontfamily\DisplayBlk{Inter Display Black}[Ligatures=TeX]
\newfontfamily\DisplaySB{Inter Display SemiBold}[Ligatures=TeX]
\newfontfamily\DisplayThin{Inter Display Thin}[Ligatures=TeX]

% --- COLOR (mismas 4 paletas que el libro) ----------------------------
\definecolor{mat} {RGB}{ 30,  64, 120}
\definecolor{cnat}{RGB}{ 16, 110,  80}
\definecolor{csoc}{RGB}{165,  75,  45}
\definecolor{eco} {RGB}{ 70,  55, 130}
\definecolor{ink}{RGB}{28,28,30}        % casi-negro para texto
\definecolor{paper}{RGB}{250,250,248}   % blanco cálido de fondo
\colorlet{subjectcolor}{csoc}
% \setsubject re-aplica los colores de beamer además de actualizar el alias
\newcommand{\setsubject}[1]{%
  \colorlet{subjectcolor}{#1}%
  \setbeamercolor{structure}{fg=#1}%
  \setbeamercolor{frametitle}{fg=#1,bg=}%
}

\setbeamercolor{normal text}{fg=ink, bg=paper}
\setbeamercolor{structure}{fg=subjectcolor}
\setbeamerfont{frametitle}{family=\DisplaySB}
\setbeamercolor{frametitle}{fg=subjectcolor, bg=}
\setbeamertemplate{frametitle}{%
  \vspace{0.6em}\hspace{0.2em}%
  {\DisplaySB\large\color{subjectcolor}\insertframetitle}\par
  \vspace{-0.2em}{\color{subjectcolor!30}\rule{\dimexpr\paperwidth-1.4em\relax}{1.2pt}}}

% Pie minimal: instituto · profe · tema · número en cuadro de color
\newcommand{\bmrtema}{}
\newcommand{\bmrprofe}{}   % ← \renewcommand{\bmrprofe}{Prof. Nombre} en cada beamer
\setbeamertemplate{footline}{%
  \hbox{}\vspace{2pt}\hspace{0.6em}%
  {\footnotesize\color{ink!45}ITSTZ%
   \ifx\bmrprofe\empty\else\space·\space\bmrprofe\fi}%
  \hfill{\footnotesize\color{ink!45}\bmrtema}\hspace{0.6em}%
  \tikz[baseline=-3pt]{\node[fill=subjectcolor, text=white, inner xsep=6pt,
     inner ysep=2.5pt, font=\DisplaySB\scriptsize]{\insertframenumber};}%
  \hspace{0.6em}\vspace{4pt}}

% =====================================================================
%  PORTADA y DIVISORES
% =====================================================================
\newcommand{\portada}[3]{% {kicker}{TÍTULO}{subtítulo}
{\setbeamertemplate{footline}{}
\begin{frame}[plain]
\begin{tikzpicture}[remember picture, overlay]
  \fill[paper] (current page.south west) rectangle (current page.north east);
  \node[anchor=north west, text=subjectcolor!12, font=\DisplayBlk]
    at ([xshift=-0.5cm,yshift=1.2cm]current page.north west)
    {\fontsize{190}{190}\selectfont ·};
  \node[anchor=west, text=subjectcolor, font=\DisplaySB\Large]
    at ([xshift=1.2cm, yshift=1.6cm]current page.west) {#1};
  \node[anchor=west, text=ink, font=\DisplayBlk, align=left]
    at ([xshift=1.15cm, yshift=0.2cm]current page.west)
    {\fontsize{40}{44}\selectfont #2};
  \fill[subjectcolor] ([xshift=1.2cm, yshift=-1.4cm]current page.west)
    rectangle ++(3.2cm, 0.16cm);
  \node[anchor=west, text=ink!60, font=\Display, align=left, text width=0.7\paperwidth]
    at ([xshift=1.2cm, yshift=-2.3cm]current page.west) {#3};
\end{tikzpicture}
\end{frame}}}

% Lámina divisoria: solo una frase grande centrada sobre color
\newcommand{\seccion}[1]{%
{\setbeamertemplate{footline}{}
\begin{frame}[plain]
\begin{tikzpicture}[remember picture, overlay]
  \fill[subjectcolor] (current page.south west) rectangle (current page.north east);
  \node[text=white, font=\DisplaySB, align=center, text width=0.85\paperwidth]
    at (current page.center) {\fontsize{30}{36}\selectfont #1};
\end{tikzpicture}
\end{frame}}}

% =====================================================================
%  ELEMENTOS DE CONTENIDO
% =====================================================================
% Una afirmación gigante que llena la lámina
\newcommand{\idea}[1]{%
{\setbeamertemplate{footline}{}
\begin{frame}[plain]
\begin{tikzpicture}[remember picture, overlay]
  \fill[paper] (current page.south west) rectangle (current page.north east);
  \fill[subjectcolor] (current page.north west)
    rectangle ([yshift=-0.32cm]current page.north east);
  \node[text=ink, font=\DisplaySB, align=center,
        text width=0.82\paperwidth]
    at (current page.center)
    {\fontsize{32}{44}\selectfont #1};
\end{tikzpicture}
\end{frame}}}

% Realce inline (fondo color suave, texto oscuro)
\newcommand{\resalta}[1]{\tikz[baseline=(R.base)]{\node[fill=subjectcolor!18,
  rounded corners=2pt, inner xsep=3pt, inner ysep=1pt](R){#1};}}
% Realce fuerte (fondo sólido, texto blanco)
\newcommand{\marca}[1]{\tikz[baseline=(R.base)]{\node[fill=subjectcolor,
  text=white, rounded corners=2pt, inner xsep=4pt, inner ysep=1.5pt](R){\textbf{#1}};}}

% Pregunta viva de apertura
\newcommand{\pregunta}[1]{%
\begin{frame}[plain]
\begin{tikzpicture}[remember picture,overlay]
  \fill[subjectcolor!7] (current page.south west) rectangle (current page.north east);
  \fill[subjectcolor] (current page.south west)
    rectangle ([xshift=0.45cm]current page.north west);
\end{tikzpicture}
\vfill
\hspace{1.5cm}\begin{minipage}{0.78\paperwidth}
  {\color{subjectcolor}\DisplayBlk\fontsize{54}{54}\selectfont ¿}\par
  \vspace{0.25em}
  {\Display\color{ink}\fontsize{22}{28}\selectfont #1}\par
  \vspace{0.4em}
  {\hfill\color{subjectcolor}\DisplayBlk\fontsize{38}{38}\selectfont ?}
\end{minipage}
\vfill
\end{frame}}

% Cierre / acción — label opcional (default: "DE LA QUEJA A LA ACCIÓN")
% Uso: \accion{texto}  o  \accion[ETIQUETA CUSTOM]{texto}
\newcommand{\accion}[2][DE LA QUEJA A LA ACCIÓN]{%
\begin{frame}[plain]
\begin{tikzpicture}[remember picture, overlay]
  \fill[subjectcolor!10] (current page.south west) rectangle (current page.north east);
  \node[anchor=west, text=subjectcolor, font=\DisplaySB]
    at ([xshift=1cm, yshift=2.2cm]current page.west)
    {\large #1};
  \node[anchor=west, text=ink, font=\Display, align=left, text width=0.78\paperwidth]
    at ([xshift=1cm, yshift=0cm]current page.west)
    {\fontsize{20}{26}\selectfont #2};
\end{tikzpicture}
\end{frame}}

% =====================================================================
%  LÍNEA DE TIEMPO HORIZONTAL
%  \begin{lineah}{año_ini}{año_fin}
%     \hito{posición 0-1}{etiqueta}   ...
%  \end{lineah}
% =====================================================================
\newenvironment{lineah}{%
  \begin{center}\begin{tikzpicture}[x=0.92\paperwidth, font=\footnotesize\sffamily]
  \draw[line width=2pt, subjectcolor, -{Stealth[length=5pt]}] (0,0) -- (1.02,0);
}{%
  \end{tikzpicture}\end{center}}
% hito: #1 = posición relativa 0..1 ; #2 = etiqueta (texto corto)
\newcommand{\hito}[2]{%
  \fill[subjectcolor] (#1,0) circle (3.2pt);
  \node[above, text=ink, align=center, text width=2.6cm, anchor=south,
        yshift=2pt] at (#1,0) {\bfseries #2};}
% hito bajo la línea (para alternar y no chocar)
\newcommand{\hitob}[2]{%
  \fill[subjectcolor] (#1,0) circle (3.2pt);
  \node[below, text=ink, align=center, text width=2.6cm, anchor=north,
        yshift=-2pt] at (#1,0) {#2};}

% =====================================================================
%  LÍNEA DE TIEMPO VERTICAL (ideal para listas cronológicas largas)
%  \begin{lineav} \hitov{año}{Título}{detalle} ... \end{lineav}
% =====================================================================
\newcounter{lvrow}
\newenvironment{lineav}{%
  \setcounter{lvrow}{0}%
  \begin{list}{}{\leftmargin=1.6cm \itemsep=0.55em \topsep=0.3em}%
  \renewcommand{\hitov}[3]{\stepcounter{lvrow}%
    \item[\tikz{\node[circle, fill=subjectcolor, text=white, inner sep=0pt,
       minimum size=1.5em, font=\DisplaySB\scriptsize]{\thelvrow};}]%
    {\DisplaySB\color{subjectcolor}##1}\quad{\bfseries ##2}\\[1pt]
    {\small\color{ink!75}##3}}%
}{%
  \end{list}}
\newcommand{\hitov}[3]{}  % placeholder; redefinido dentro del entorno

% =====================================================================
%  MAPA MENTAL — nodo central + ramas radiales
%  \begin{mapa}{NODO CENTRAL}
%     \rama{ángulo}{distancia cm}{texto de la rama}  ...
%  \end{mapa}
% =====================================================================
\newenvironment{mapa}[2][1]{%   % [escala opcional, def=1]{NODO CENTRAL}
  \begin{center}\begin{tikzpicture}[font=\sffamily\small, scale=#1, transform shape]
  \node[circle, fill=subjectcolor, text=white, font=\DisplaySB,
        align=center, text width=2.4cm, inner sep=4pt, minimum size=2.6cm]
        (CENTRO) at (0,0) {#2};
  \renewcommand{\rama}[3]{%
    \node[draw=subjectcolor!55, fill=subjectcolor!8, rounded corners=4pt,
          text=ink, align=center, text width=2.7cm, inner sep=4pt]
          (R##1) at (##1:##2) {##3};
    \draw[subjectcolor!55, line width=1.2pt] (CENTRO) -- (R##1);}%
}{%
  \end{tikzpicture}\end{center}}
\newcommand{\rama}[3]{}  % placeholder; redefinido dentro del entorno
 ; \setsubject{csoc} ; \begin{document}...
% =====================================================================
\documentclass[aspectratio=169, 11pt]{beamer}
\usetheme{default}
\setbeamertemplate{navigation symbols}{}   % sin botones de navegación

\usepackage{fontspec}
\usepackage{polyglossia}\setdefaultlanguage{spanish}
\usepackage{tikz}
\usetikzlibrary{calc, positioning, arrows.meta, decorations.pathreplacing}
\usepackage{amsmath}

% --- TIPOGRAFÍA (idéntica familia que el libro) -----------------------
\setsansfont{Inter}[UprightFont=Inter-Regular, BoldFont=Inter-SemiBold,
  ItalicFont=Inter-Italic, Ligatures=TeX, Numbers=Lining]
\newfontfamily\Display{Inter Display}[Ligatures=TeX]
\newfontfamily\DisplayBlk{Inter Display Black}[Ligatures=TeX]
\newfontfamily\DisplaySB{Inter Display SemiBold}[Ligatures=TeX]
\newfontfamily\DisplayThin{Inter Display Thin}[Ligatures=TeX]

% --- COLOR (mismas 4 paletas que el libro) ----------------------------
\definecolor{mat} {RGB}{ 30,  64, 120}
\definecolor{cnat}{RGB}{ 16, 110,  80}
\definecolor{csoc}{RGB}{165,  75,  45}
\definecolor{eco} {RGB}{ 70,  55, 130}
\definecolor{ink}{RGB}{28,28,30}        % casi-negro para texto
\definecolor{paper}{RGB}{250,250,248}   % blanco cálido de fondo
\colorlet{subjectcolor}{csoc}
% \setsubject re-aplica los colores de beamer además de actualizar el alias
\newcommand{\setsubject}[1]{%
  \colorlet{subjectcolor}{#1}%
  \setbeamercolor{structure}{fg=#1}%
  \setbeamercolor{frametitle}{fg=#1,bg=}%
}

\setbeamercolor{normal text}{fg=ink, bg=paper}
\setbeamercolor{structure}{fg=subjectcolor}
\setbeamerfont{frametitle}{family=\DisplaySB}
\setbeamercolor{frametitle}{fg=subjectcolor, bg=}
\setbeamertemplate{frametitle}{%
  \vspace{0.6em}\hspace{0.2em}%
  {\DisplaySB\large\color{subjectcolor}\insertframetitle}\par
  \vspace{-0.2em}{\color{subjectcolor!30}\rule{\dimexpr\paperwidth-1.4em\relax}{1.2pt}}}

% Pie minimal: instituto · profe · tema · número en cuadro de color
\newcommand{\bmrtema}{}
\newcommand{\bmrprofe}{}   % ← \renewcommand{\bmrprofe}{Prof. Nombre} en cada beamer
\setbeamertemplate{footline}{%
  \hbox{}\vspace{2pt}\hspace{0.6em}%
  {\footnotesize\color{ink!45}ITSTZ%
   \ifx\bmrprofe\empty\else\space·\space\bmrprofe\fi}%
  \hfill{\footnotesize\color{ink!45}\bmrtema}\hspace{0.6em}%
  \tikz[baseline=-3pt]{\node[fill=subjectcolor, text=white, inner xsep=6pt,
     inner ysep=2.5pt, font=\DisplaySB\scriptsize]{\insertframenumber};}%
  \hspace{0.6em}\vspace{4pt}}

% =====================================================================
%  PORTADA y DIVISORES
% =====================================================================
\newcommand{\portada}[3]{% {kicker}{TÍTULO}{subtítulo}
{\setbeamertemplate{footline}{}
\begin{frame}[plain]
\begin{tikzpicture}[remember picture, overlay]
  \fill[paper] (current page.south west) rectangle (current page.north east);
  \node[anchor=north west, text=subjectcolor!12, font=\DisplayBlk]
    at ([xshift=-0.5cm,yshift=1.2cm]current page.north west)
    {\fontsize{190}{190}\selectfont ·};
  \node[anchor=west, text=subjectcolor, font=\DisplaySB\Large]
    at ([xshift=1.2cm, yshift=1.6cm]current page.west) {#1};
  \node[anchor=west, text=ink, font=\DisplayBlk, align=left]
    at ([xshift=1.15cm, yshift=0.2cm]current page.west)
    {\fontsize{40}{44}\selectfont #2};
  \fill[subjectcolor] ([xshift=1.2cm, yshift=-1.4cm]current page.west)
    rectangle ++(3.2cm, 0.16cm);
  \node[anchor=west, text=ink!60, font=\Display, align=left, text width=0.7\paperwidth]
    at ([xshift=1.2cm, yshift=-2.3cm]current page.west) {#3};
\end{tikzpicture}
\end{frame}}}

% Lámina divisoria: solo una frase grande centrada sobre color
\newcommand{\seccion}[1]{%
{\setbeamertemplate{footline}{}
\begin{frame}[plain]
\begin{tikzpicture}[remember picture, overlay]
  \fill[subjectcolor] (current page.south west) rectangle (current page.north east);
  \node[text=white, font=\DisplaySB, align=center, text width=0.85\paperwidth]
    at (current page.center) {\fontsize{30}{36}\selectfont #1};
\end{tikzpicture}
\end{frame}}}

% =====================================================================
%  ELEMENTOS DE CONTENIDO
% =====================================================================
% Una afirmación gigante que llena la lámina
\newcommand{\idea}[1]{%
{\setbeamertemplate{footline}{}
\begin{frame}[plain]
\begin{tikzpicture}[remember picture, overlay]
  \fill[paper] (current page.south west) rectangle (current page.north east);
  \fill[subjectcolor] (current page.north west)
    rectangle ([yshift=-0.32cm]current page.north east);
  \node[text=ink, font=\DisplaySB, align=center,
        text width=0.82\paperwidth]
    at (current page.center)
    {\fontsize{32}{44}\selectfont #1};
\end{tikzpicture}
\end{frame}}}

% Realce inline (fondo color suave, texto oscuro)
\newcommand{\resalta}[1]{\tikz[baseline=(R.base)]{\node[fill=subjectcolor!18,
  rounded corners=2pt, inner xsep=3pt, inner ysep=1pt](R){#1};}}
% Realce fuerte (fondo sólido, texto blanco)
\newcommand{\marca}[1]{\tikz[baseline=(R.base)]{\node[fill=subjectcolor,
  text=white, rounded corners=2pt, inner xsep=4pt, inner ysep=1.5pt](R){\textbf{#1}};}}

% Pregunta viva de apertura
\newcommand{\pregunta}[1]{%
\begin{frame}[plain]
\begin{tikzpicture}[remember picture,overlay]
  \fill[subjectcolor!7] (current page.south west) rectangle (current page.north east);
  \fill[subjectcolor] (current page.south west)
    rectangle ([xshift=0.45cm]current page.north west);
\end{tikzpicture}
\vfill
\hspace{1.5cm}\begin{minipage}{0.78\paperwidth}
  {\color{subjectcolor}\DisplayBlk\fontsize{54}{54}\selectfont ¿}\par
  \vspace{0.25em}
  {\Display\color{ink}\fontsize{22}{28}\selectfont #1}\par
  \vspace{0.4em}
  {\hfill\color{subjectcolor}\DisplayBlk\fontsize{38}{38}\selectfont ?}
\end{minipage}
\vfill
\end{frame}}

% Cierre / acción — label opcional (default: "DE LA QUEJA A LA ACCIÓN")
% Uso: \accion{texto}  o  \accion[ETIQUETA CUSTOM]{texto}
\newcommand{\accion}[2][DE LA QUEJA A LA ACCIÓN]{%
\begin{frame}[plain]
\begin{tikzpicture}[remember picture, overlay]
  \fill[subjectcolor!10] (current page.south west) rectangle (current page.north east);
  \node[anchor=west, text=subjectcolor, font=\DisplaySB]
    at ([xshift=1cm, yshift=2.2cm]current page.west)
    {\large #1};
  \node[anchor=west, text=ink, font=\Display, align=left, text width=0.78\paperwidth]
    at ([xshift=1cm, yshift=0cm]current page.west)
    {\fontsize{20}{26}\selectfont #2};
\end{tikzpicture}
\end{frame}}

% =====================================================================
%  LÍNEA DE TIEMPO HORIZONTAL
%  \begin{lineah}{año_ini}{año_fin}
%     \hito{posición 0-1}{etiqueta}   ...
%  \end{lineah}
% =====================================================================
\newenvironment{lineah}{%
  \begin{center}\begin{tikzpicture}[x=0.92\paperwidth, font=\footnotesize\sffamily]
  \draw[line width=2pt, subjectcolor, -{Stealth[length=5pt]}] (0,0) -- (1.02,0);
}{%
  \end{tikzpicture}\end{center}}
% hito: #1 = posición relativa 0..1 ; #2 = etiqueta (texto corto)
\newcommand{\hito}[2]{%
  \fill[subjectcolor] (#1,0) circle (3.2pt);
  \node[above, text=ink, align=center, text width=2.6cm, anchor=south,
        yshift=2pt] at (#1,0) {\bfseries #2};}
% hito bajo la línea (para alternar y no chocar)
\newcommand{\hitob}[2]{%
  \fill[subjectcolor] (#1,0) circle (3.2pt);
  \node[below, text=ink, align=center, text width=2.6cm, anchor=north,
        yshift=-2pt] at (#1,0) {#2};}

% =====================================================================
%  LÍNEA DE TIEMPO VERTICAL (ideal para listas cronológicas largas)
%  \begin{lineav} \hitov{año}{Título}{detalle} ... \end{lineav}
% =====================================================================
\newcounter{lvrow}
\newenvironment{lineav}{%
  \setcounter{lvrow}{0}%
  \begin{list}{}{\leftmargin=1.6cm \itemsep=0.55em \topsep=0.3em}%
  \renewcommand{\hitov}[3]{\stepcounter{lvrow}%
    \item[\tikz{\node[circle, fill=subjectcolor, text=white, inner sep=0pt,
       minimum size=1.5em, font=\DisplaySB\scriptsize]{\thelvrow};}]%
    {\DisplaySB\color{subjectcolor}##1}\quad{\bfseries ##2}\\[1pt]
    {\small\color{ink!75}##3}}%
}{%
  \end{list}}
\newcommand{\hitov}[3]{}  % placeholder; redefinido dentro del entorno

% =====================================================================
%  MAPA MENTAL — nodo central + ramas radiales
%  \begin{mapa}{NODO CENTRAL}
%     \rama{ángulo}{distancia cm}{texto de la rama}  ...
%  \end{mapa}
% =====================================================================
\newenvironment{mapa}[2][1]{%   % [escala opcional, def=1]{NODO CENTRAL}
  \begin{center}\begin{tikzpicture}[font=\sffamily\small, scale=#1, transform shape]
  \node[circle, fill=subjectcolor, text=white, font=\DisplaySB,
        align=center, text width=2.4cm, inner sep=4pt, minimum size=2.6cm]
        (CENTRO) at (0,0) {#2};
  \renewcommand{\rama}[3]{%
    \node[draw=subjectcolor!55, fill=subjectcolor!8, rounded corners=4pt,
          text=ink, align=center, text width=2.7cm, inner sep=4pt]
          (R##1) at (##1:##2) {##3};
    \draw[subjectcolor!55, line width=1.2pt] (CENTRO) -- (R##1);}%
}{%
  \end{tikzpicture}\end{center}}
\newcommand{\rama}[3]{}  % placeholder; redefinido dentro del entorno
):
%     % =====================================================================
%  DESIGN SYSTEM BEAMER v5 — ITSTZ · 16:9
%  Coherente con el libro: Inter Display (títulos) + Inter (cuerpo),
%  un color de asignatura, fondo claro, texto GRANDE y contrastado.
%  Filosofía slide: una idea por lámina, mucho aire, cero relleno.
%
%  Vocabulario:
%    \portada{kicker}{TÍTULO}{subtítulo}
%    \seccion{texto}                  lámina divisoria a pantalla completa
%    \idea{TEXTO GRANDE}              una afirmación que llena la slide
%    \resalta{texto}                  realce inline (fondo color)
%    frames de contenido normales con \frametitle
%    \begin{lineah} \hito{x}{etiqueta} ... \end{lineah}   timeline horizontal
%    \begin{lineav} \hitov{x}{título}{texto} ... \end{lineav} timeline vertical
%    \begin{mapa}{NODO CENTRAL} \rama{ang}{dist}{texto} ... \end{mapa}
%    \pregunta{texto}                 pregunta viva (apertura)
%    \accion{texto}                   cierre "de la queja a la acción"
%  Uso: % =====================================================================
%  DESIGN SYSTEM BEAMER v5 — ITSTZ · 16:9
%  Coherente con el libro: Inter Display (títulos) + Inter (cuerpo),
%  un color de asignatura, fondo claro, texto GRANDE y contrastado.
%  Filosofía slide: una idea por lámina, mucho aire, cero relleno.
%
%  Vocabulario:
%    \portada{kicker}{TÍTULO}{subtítulo}
%    \seccion{texto}                  lámina divisoria a pantalla completa
%    \idea{TEXTO GRANDE}              una afirmación que llena la slide
%    \resalta{texto}                  realce inline (fondo color)
%    frames de contenido normales con \frametitle
%    \begin{lineah} \hito{x}{etiqueta} ... \end{lineah}   timeline horizontal
%    \begin{lineav} \hitov{x}{título}{texto} ... \end{lineav} timeline vertical
%    \begin{mapa}{NODO CENTRAL} \rama{ang}{dist}{texto} ... \end{mapa}
%    \pregunta{texto}                 pregunta viva (apertura)
%    \accion{texto}                   cierre "de la queja a la acción"
%  Uso: % =====================================================================
%  DESIGN SYSTEM BEAMER v5 — ITSTZ · 16:9
%  Coherente con el libro: Inter Display (títulos) + Inter (cuerpo),
%  un color de asignatura, fondo claro, texto GRANDE y contrastado.
%  Filosofía slide: una idea por lámina, mucho aire, cero relleno.
%
%  Vocabulario:
%    \portada{kicker}{TÍTULO}{subtítulo}
%    \seccion{texto}                  lámina divisoria a pantalla completa
%    \idea{TEXTO GRANDE}              una afirmación que llena la slide
%    \resalta{texto}                  realce inline (fondo color)
%    frames de contenido normales con \frametitle
%    \begin{lineah} \hito{x}{etiqueta} ... \end{lineah}   timeline horizontal
%    \begin{lineav} \hitov{x}{título}{texto} ... \end{lineav} timeline vertical
%    \begin{mapa}{NODO CENTRAL} \rama{ang}{dist}{texto} ... \end{mapa}
%    \pregunta{texto}                 pregunta viva (apertura)
%    \accion{texto}                   cierre "de la queja a la acción"
%  Uso: \input{design/beamer.tex} ; \setsubject{csoc} ; \begin{document}...
% =====================================================================
\documentclass[aspectratio=169, 11pt]{beamer}
\usetheme{default}
\setbeamertemplate{navigation symbols}{}   % sin botones de navegación

\usepackage{fontspec}
\usepackage{polyglossia}\setdefaultlanguage{spanish}
\usepackage{tikz}
\usetikzlibrary{calc, positioning, arrows.meta, decorations.pathreplacing}
\usepackage{amsmath}

% --- TIPOGRAFÍA (idéntica familia que el libro) -----------------------
\setsansfont{Inter}[UprightFont=Inter-Regular, BoldFont=Inter-SemiBold,
  ItalicFont=Inter-Italic, Ligatures=TeX, Numbers=Lining]
\newfontfamily\Display{Inter Display}[Ligatures=TeX]
\newfontfamily\DisplayBlk{Inter Display Black}[Ligatures=TeX]
\newfontfamily\DisplaySB{Inter Display SemiBold}[Ligatures=TeX]
\newfontfamily\DisplayThin{Inter Display Thin}[Ligatures=TeX]

% --- COLOR (mismas 4 paletas que el libro) ----------------------------
\definecolor{mat} {RGB}{ 30,  64, 120}
\definecolor{cnat}{RGB}{ 16, 110,  80}
\definecolor{csoc}{RGB}{165,  75,  45}
\definecolor{eco} {RGB}{ 70,  55, 130}
\definecolor{ink}{RGB}{28,28,30}        % casi-negro para texto
\definecolor{paper}{RGB}{250,250,248}   % blanco cálido de fondo
\colorlet{subjectcolor}{csoc}
% \setsubject re-aplica los colores de beamer además de actualizar el alias
\newcommand{\setsubject}[1]{%
  \colorlet{subjectcolor}{#1}%
  \setbeamercolor{structure}{fg=#1}%
  \setbeamercolor{frametitle}{fg=#1,bg=}%
}

\setbeamercolor{normal text}{fg=ink, bg=paper}
\setbeamercolor{structure}{fg=subjectcolor}
\setbeamerfont{frametitle}{family=\DisplaySB}
\setbeamercolor{frametitle}{fg=subjectcolor, bg=}
\setbeamertemplate{frametitle}{%
  \vspace{0.6em}\hspace{0.2em}%
  {\DisplaySB\large\color{subjectcolor}\insertframetitle}\par
  \vspace{-0.2em}{\color{subjectcolor!30}\rule{\dimexpr\paperwidth-1.4em\relax}{1.2pt}}}

% Pie minimal: instituto · profe · tema · número en cuadro de color
\newcommand{\bmrtema}{}
\newcommand{\bmrprofe}{}   % ← \renewcommand{\bmrprofe}{Prof. Nombre} en cada beamer
\setbeamertemplate{footline}{%
  \hbox{}\vspace{2pt}\hspace{0.6em}%
  {\footnotesize\color{ink!45}ITSTZ%
   \ifx\bmrprofe\empty\else\space·\space\bmrprofe\fi}%
  \hfill{\footnotesize\color{ink!45}\bmrtema}\hspace{0.6em}%
  \tikz[baseline=-3pt]{\node[fill=subjectcolor, text=white, inner xsep=6pt,
     inner ysep=2.5pt, font=\DisplaySB\scriptsize]{\insertframenumber};}%
  \hspace{0.6em}\vspace{4pt}}

% =====================================================================
%  PORTADA y DIVISORES
% =====================================================================
\newcommand{\portada}[3]{% {kicker}{TÍTULO}{subtítulo}
{\setbeamertemplate{footline}{}
\begin{frame}[plain]
\begin{tikzpicture}[remember picture, overlay]
  \fill[paper] (current page.south west) rectangle (current page.north east);
  \node[anchor=north west, text=subjectcolor!12, font=\DisplayBlk]
    at ([xshift=-0.5cm,yshift=1.2cm]current page.north west)
    {\fontsize{190}{190}\selectfont ·};
  \node[anchor=west, text=subjectcolor, font=\DisplaySB\Large]
    at ([xshift=1.2cm, yshift=1.6cm]current page.west) {#1};
  \node[anchor=west, text=ink, font=\DisplayBlk, align=left]
    at ([xshift=1.15cm, yshift=0.2cm]current page.west)
    {\fontsize{40}{44}\selectfont #2};
  \fill[subjectcolor] ([xshift=1.2cm, yshift=-1.4cm]current page.west)
    rectangle ++(3.2cm, 0.16cm);
  \node[anchor=west, text=ink!60, font=\Display, align=left, text width=0.7\paperwidth]
    at ([xshift=1.2cm, yshift=-2.3cm]current page.west) {#3};
\end{tikzpicture}
\end{frame}}}

% Lámina divisoria: solo una frase grande centrada sobre color
\newcommand{\seccion}[1]{%
{\setbeamertemplate{footline}{}
\begin{frame}[plain]
\begin{tikzpicture}[remember picture, overlay]
  \fill[subjectcolor] (current page.south west) rectangle (current page.north east);
  \node[text=white, font=\DisplaySB, align=center, text width=0.85\paperwidth]
    at (current page.center) {\fontsize{30}{36}\selectfont #1};
\end{tikzpicture}
\end{frame}}}

% =====================================================================
%  ELEMENTOS DE CONTENIDO
% =====================================================================
% Una afirmación gigante que llena la lámina
\newcommand{\idea}[1]{%
{\setbeamertemplate{footline}{}
\begin{frame}[plain]
\begin{tikzpicture}[remember picture, overlay]
  \fill[paper] (current page.south west) rectangle (current page.north east);
  \fill[subjectcolor] (current page.north west)
    rectangle ([yshift=-0.32cm]current page.north east);
  \node[text=ink, font=\DisplaySB, align=center,
        text width=0.82\paperwidth]
    at (current page.center)
    {\fontsize{32}{44}\selectfont #1};
\end{tikzpicture}
\end{frame}}}

% Realce inline (fondo color suave, texto oscuro)
\newcommand{\resalta}[1]{\tikz[baseline=(R.base)]{\node[fill=subjectcolor!18,
  rounded corners=2pt, inner xsep=3pt, inner ysep=1pt](R){#1};}}
% Realce fuerte (fondo sólido, texto blanco)
\newcommand{\marca}[1]{\tikz[baseline=(R.base)]{\node[fill=subjectcolor,
  text=white, rounded corners=2pt, inner xsep=4pt, inner ysep=1.5pt](R){\textbf{#1}};}}

% Pregunta viva de apertura
\newcommand{\pregunta}[1]{%
\begin{frame}[plain]
\begin{tikzpicture}[remember picture,overlay]
  \fill[subjectcolor!7] (current page.south west) rectangle (current page.north east);
  \fill[subjectcolor] (current page.south west)
    rectangle ([xshift=0.45cm]current page.north west);
\end{tikzpicture}
\vfill
\hspace{1.5cm}\begin{minipage}{0.78\paperwidth}
  {\color{subjectcolor}\DisplayBlk\fontsize{54}{54}\selectfont ¿}\par
  \vspace{0.25em}
  {\Display\color{ink}\fontsize{22}{28}\selectfont #1}\par
  \vspace{0.4em}
  {\hfill\color{subjectcolor}\DisplayBlk\fontsize{38}{38}\selectfont ?}
\end{minipage}
\vfill
\end{frame}}

% Cierre / acción — label opcional (default: "DE LA QUEJA A LA ACCIÓN")
% Uso: \accion{texto}  o  \accion[ETIQUETA CUSTOM]{texto}
\newcommand{\accion}[2][DE LA QUEJA A LA ACCIÓN]{%
\begin{frame}[plain]
\begin{tikzpicture}[remember picture, overlay]
  \fill[subjectcolor!10] (current page.south west) rectangle (current page.north east);
  \node[anchor=west, text=subjectcolor, font=\DisplaySB]
    at ([xshift=1cm, yshift=2.2cm]current page.west)
    {\large #1};
  \node[anchor=west, text=ink, font=\Display, align=left, text width=0.78\paperwidth]
    at ([xshift=1cm, yshift=0cm]current page.west)
    {\fontsize{20}{26}\selectfont #2};
\end{tikzpicture}
\end{frame}}

% =====================================================================
%  LÍNEA DE TIEMPO HORIZONTAL
%  \begin{lineah}{año_ini}{año_fin}
%     \hito{posición 0-1}{etiqueta}   ...
%  \end{lineah}
% =====================================================================
\newenvironment{lineah}{%
  \begin{center}\begin{tikzpicture}[x=0.92\paperwidth, font=\footnotesize\sffamily]
  \draw[line width=2pt, subjectcolor, -{Stealth[length=5pt]}] (0,0) -- (1.02,0);
}{%
  \end{tikzpicture}\end{center}}
% hito: #1 = posición relativa 0..1 ; #2 = etiqueta (texto corto)
\newcommand{\hito}[2]{%
  \fill[subjectcolor] (#1,0) circle (3.2pt);
  \node[above, text=ink, align=center, text width=2.6cm, anchor=south,
        yshift=2pt] at (#1,0) {\bfseries #2};}
% hito bajo la línea (para alternar y no chocar)
\newcommand{\hitob}[2]{%
  \fill[subjectcolor] (#1,0) circle (3.2pt);
  \node[below, text=ink, align=center, text width=2.6cm, anchor=north,
        yshift=-2pt] at (#1,0) {#2};}

% =====================================================================
%  LÍNEA DE TIEMPO VERTICAL (ideal para listas cronológicas largas)
%  \begin{lineav} \hitov{año}{Título}{detalle} ... \end{lineav}
% =====================================================================
\newcounter{lvrow}
\newenvironment{lineav}{%
  \setcounter{lvrow}{0}%
  \begin{list}{}{\leftmargin=1.6cm \itemsep=0.55em \topsep=0.3em}%
  \renewcommand{\hitov}[3]{\stepcounter{lvrow}%
    \item[\tikz{\node[circle, fill=subjectcolor, text=white, inner sep=0pt,
       minimum size=1.5em, font=\DisplaySB\scriptsize]{\thelvrow};}]%
    {\DisplaySB\color{subjectcolor}##1}\quad{\bfseries ##2}\\[1pt]
    {\small\color{ink!75}##3}}%
}{%
  \end{list}}
\newcommand{\hitov}[3]{}  % placeholder; redefinido dentro del entorno

% =====================================================================
%  MAPA MENTAL — nodo central + ramas radiales
%  \begin{mapa}{NODO CENTRAL}
%     \rama{ángulo}{distancia cm}{texto de la rama}  ...
%  \end{mapa}
% =====================================================================
\newenvironment{mapa}[2][1]{%   % [escala opcional, def=1]{NODO CENTRAL}
  \begin{center}\begin{tikzpicture}[font=\sffamily\small, scale=#1, transform shape]
  \node[circle, fill=subjectcolor, text=white, font=\DisplaySB,
        align=center, text width=2.4cm, inner sep=4pt, minimum size=2.6cm]
        (CENTRO) at (0,0) {#2};
  \renewcommand{\rama}[3]{%
    \node[draw=subjectcolor!55, fill=subjectcolor!8, rounded corners=4pt,
          text=ink, align=center, text width=2.7cm, inner sep=4pt]
          (R##1) at (##1:##2) {##3};
    \draw[subjectcolor!55, line width=1.2pt] (CENTRO) -- (R##1);}%
}{%
  \end{tikzpicture}\end{center}}
\newcommand{\rama}[3]{}  % placeholder; redefinido dentro del entorno
 ; \setsubject{csoc} ; \begin{document}...
% =====================================================================
\documentclass[aspectratio=169, 11pt]{beamer}
\usetheme{default}
\setbeamertemplate{navigation symbols}{}   % sin botones de navegación

\usepackage{fontspec}
\usepackage{polyglossia}\setdefaultlanguage{spanish}
\usepackage{tikz}
\usetikzlibrary{calc, positioning, arrows.meta, decorations.pathreplacing}
\usepackage{amsmath}

% --- TIPOGRAFÍA (idéntica familia que el libro) -----------------------
\setsansfont{Inter}[UprightFont=Inter-Regular, BoldFont=Inter-SemiBold,
  ItalicFont=Inter-Italic, Ligatures=TeX, Numbers=Lining]
\newfontfamily\Display{Inter Display}[Ligatures=TeX]
\newfontfamily\DisplayBlk{Inter Display Black}[Ligatures=TeX]
\newfontfamily\DisplaySB{Inter Display SemiBold}[Ligatures=TeX]
\newfontfamily\DisplayThin{Inter Display Thin}[Ligatures=TeX]

% --- COLOR (mismas 4 paletas que el libro) ----------------------------
\definecolor{mat} {RGB}{ 30,  64, 120}
\definecolor{cnat}{RGB}{ 16, 110,  80}
\definecolor{csoc}{RGB}{165,  75,  45}
\definecolor{eco} {RGB}{ 70,  55, 130}
\definecolor{ink}{RGB}{28,28,30}        % casi-negro para texto
\definecolor{paper}{RGB}{250,250,248}   % blanco cálido de fondo
\colorlet{subjectcolor}{csoc}
% \setsubject re-aplica los colores de beamer además de actualizar el alias
\newcommand{\setsubject}[1]{%
  \colorlet{subjectcolor}{#1}%
  \setbeamercolor{structure}{fg=#1}%
  \setbeamercolor{frametitle}{fg=#1,bg=}%
}

\setbeamercolor{normal text}{fg=ink, bg=paper}
\setbeamercolor{structure}{fg=subjectcolor}
\setbeamerfont{frametitle}{family=\DisplaySB}
\setbeamercolor{frametitle}{fg=subjectcolor, bg=}
\setbeamertemplate{frametitle}{%
  \vspace{0.6em}\hspace{0.2em}%
  {\DisplaySB\large\color{subjectcolor}\insertframetitle}\par
  \vspace{-0.2em}{\color{subjectcolor!30}\rule{\dimexpr\paperwidth-1.4em\relax}{1.2pt}}}

% Pie minimal: instituto · profe · tema · número en cuadro de color
\newcommand{\bmrtema}{}
\newcommand{\bmrprofe}{}   % ← \renewcommand{\bmrprofe}{Prof. Nombre} en cada beamer
\setbeamertemplate{footline}{%
  \hbox{}\vspace{2pt}\hspace{0.6em}%
  {\footnotesize\color{ink!45}ITSTZ%
   \ifx\bmrprofe\empty\else\space·\space\bmrprofe\fi}%
  \hfill{\footnotesize\color{ink!45}\bmrtema}\hspace{0.6em}%
  \tikz[baseline=-3pt]{\node[fill=subjectcolor, text=white, inner xsep=6pt,
     inner ysep=2.5pt, font=\DisplaySB\scriptsize]{\insertframenumber};}%
  \hspace{0.6em}\vspace{4pt}}

% =====================================================================
%  PORTADA y DIVISORES
% =====================================================================
\newcommand{\portada}[3]{% {kicker}{TÍTULO}{subtítulo}
{\setbeamertemplate{footline}{}
\begin{frame}[plain]
\begin{tikzpicture}[remember picture, overlay]
  \fill[paper] (current page.south west) rectangle (current page.north east);
  \node[anchor=north west, text=subjectcolor!12, font=\DisplayBlk]
    at ([xshift=-0.5cm,yshift=1.2cm]current page.north west)
    {\fontsize{190}{190}\selectfont ·};
  \node[anchor=west, text=subjectcolor, font=\DisplaySB\Large]
    at ([xshift=1.2cm, yshift=1.6cm]current page.west) {#1};
  \node[anchor=west, text=ink, font=\DisplayBlk, align=left]
    at ([xshift=1.15cm, yshift=0.2cm]current page.west)
    {\fontsize{40}{44}\selectfont #2};
  \fill[subjectcolor] ([xshift=1.2cm, yshift=-1.4cm]current page.west)
    rectangle ++(3.2cm, 0.16cm);
  \node[anchor=west, text=ink!60, font=\Display, align=left, text width=0.7\paperwidth]
    at ([xshift=1.2cm, yshift=-2.3cm]current page.west) {#3};
\end{tikzpicture}
\end{frame}}}

% Lámina divisoria: solo una frase grande centrada sobre color
\newcommand{\seccion}[1]{%
{\setbeamertemplate{footline}{}
\begin{frame}[plain]
\begin{tikzpicture}[remember picture, overlay]
  \fill[subjectcolor] (current page.south west) rectangle (current page.north east);
  \node[text=white, font=\DisplaySB, align=center, text width=0.85\paperwidth]
    at (current page.center) {\fontsize{30}{36}\selectfont #1};
\end{tikzpicture}
\end{frame}}}

% =====================================================================
%  ELEMENTOS DE CONTENIDO
% =====================================================================
% Una afirmación gigante que llena la lámina
\newcommand{\idea}[1]{%
{\setbeamertemplate{footline}{}
\begin{frame}[plain]
\begin{tikzpicture}[remember picture, overlay]
  \fill[paper] (current page.south west) rectangle (current page.north east);
  \fill[subjectcolor] (current page.north west)
    rectangle ([yshift=-0.32cm]current page.north east);
  \node[text=ink, font=\DisplaySB, align=center,
        text width=0.82\paperwidth]
    at (current page.center)
    {\fontsize{32}{44}\selectfont #1};
\end{tikzpicture}
\end{frame}}}

% Realce inline (fondo color suave, texto oscuro)
\newcommand{\resalta}[1]{\tikz[baseline=(R.base)]{\node[fill=subjectcolor!18,
  rounded corners=2pt, inner xsep=3pt, inner ysep=1pt](R){#1};}}
% Realce fuerte (fondo sólido, texto blanco)
\newcommand{\marca}[1]{\tikz[baseline=(R.base)]{\node[fill=subjectcolor,
  text=white, rounded corners=2pt, inner xsep=4pt, inner ysep=1.5pt](R){\textbf{#1}};}}

% Pregunta viva de apertura
\newcommand{\pregunta}[1]{%
\begin{frame}[plain]
\begin{tikzpicture}[remember picture,overlay]
  \fill[subjectcolor!7] (current page.south west) rectangle (current page.north east);
  \fill[subjectcolor] (current page.south west)
    rectangle ([xshift=0.45cm]current page.north west);
\end{tikzpicture}
\vfill
\hspace{1.5cm}\begin{minipage}{0.78\paperwidth}
  {\color{subjectcolor}\DisplayBlk\fontsize{54}{54}\selectfont ¿}\par
  \vspace{0.25em}
  {\Display\color{ink}\fontsize{22}{28}\selectfont #1}\par
  \vspace{0.4em}
  {\hfill\color{subjectcolor}\DisplayBlk\fontsize{38}{38}\selectfont ?}
\end{minipage}
\vfill
\end{frame}}

% Cierre / acción — label opcional (default: "DE LA QUEJA A LA ACCIÓN")
% Uso: \accion{texto}  o  \accion[ETIQUETA CUSTOM]{texto}
\newcommand{\accion}[2][DE LA QUEJA A LA ACCIÓN]{%
\begin{frame}[plain]
\begin{tikzpicture}[remember picture, overlay]
  \fill[subjectcolor!10] (current page.south west) rectangle (current page.north east);
  \node[anchor=west, text=subjectcolor, font=\DisplaySB]
    at ([xshift=1cm, yshift=2.2cm]current page.west)
    {\large #1};
  \node[anchor=west, text=ink, font=\Display, align=left, text width=0.78\paperwidth]
    at ([xshift=1cm, yshift=0cm]current page.west)
    {\fontsize{20}{26}\selectfont #2};
\end{tikzpicture}
\end{frame}}

% =====================================================================
%  LÍNEA DE TIEMPO HORIZONTAL
%  \begin{lineah}{año_ini}{año_fin}
%     \hito{posición 0-1}{etiqueta}   ...
%  \end{lineah}
% =====================================================================
\newenvironment{lineah}{%
  \begin{center}\begin{tikzpicture}[x=0.92\paperwidth, font=\footnotesize\sffamily]
  \draw[line width=2pt, subjectcolor, -{Stealth[length=5pt]}] (0,0) -- (1.02,0);
}{%
  \end{tikzpicture}\end{center}}
% hito: #1 = posición relativa 0..1 ; #2 = etiqueta (texto corto)
\newcommand{\hito}[2]{%
  \fill[subjectcolor] (#1,0) circle (3.2pt);
  \node[above, text=ink, align=center, text width=2.6cm, anchor=south,
        yshift=2pt] at (#1,0) {\bfseries #2};}
% hito bajo la línea (para alternar y no chocar)
\newcommand{\hitob}[2]{%
  \fill[subjectcolor] (#1,0) circle (3.2pt);
  \node[below, text=ink, align=center, text width=2.6cm, anchor=north,
        yshift=-2pt] at (#1,0) {#2};}

% =====================================================================
%  LÍNEA DE TIEMPO VERTICAL (ideal para listas cronológicas largas)
%  \begin{lineav} \hitov{año}{Título}{detalle} ... \end{lineav}
% =====================================================================
\newcounter{lvrow}
\newenvironment{lineav}{%
  \setcounter{lvrow}{0}%
  \begin{list}{}{\leftmargin=1.6cm \itemsep=0.55em \topsep=0.3em}%
  \renewcommand{\hitov}[3]{\stepcounter{lvrow}%
    \item[\tikz{\node[circle, fill=subjectcolor, text=white, inner sep=0pt,
       minimum size=1.5em, font=\DisplaySB\scriptsize]{\thelvrow};}]%
    {\DisplaySB\color{subjectcolor}##1}\quad{\bfseries ##2}\\[1pt]
    {\small\color{ink!75}##3}}%
}{%
  \end{list}}
\newcommand{\hitov}[3]{}  % placeholder; redefinido dentro del entorno

% =====================================================================
%  MAPA MENTAL — nodo central + ramas radiales
%  \begin{mapa}{NODO CENTRAL}
%     \rama{ángulo}{distancia cm}{texto de la rama}  ...
%  \end{mapa}
% =====================================================================
\newenvironment{mapa}[2][1]{%   % [escala opcional, def=1]{NODO CENTRAL}
  \begin{center}\begin{tikzpicture}[font=\sffamily\small, scale=#1, transform shape]
  \node[circle, fill=subjectcolor, text=white, font=\DisplaySB,
        align=center, text width=2.4cm, inner sep=4pt, minimum size=2.6cm]
        (CENTRO) at (0,0) {#2};
  \renewcommand{\rama}[3]{%
    \node[draw=subjectcolor!55, fill=subjectcolor!8, rounded corners=4pt,
          text=ink, align=center, text width=2.7cm, inner sep=4pt]
          (R##1) at (##1:##2) {##3};
    \draw[subjectcolor!55, line width=1.2pt] (CENTRO) -- (R##1);}%
}{%
  \end{tikzpicture}\end{center}}
\newcommand{\rama}[3]{}  % placeholder; redefinido dentro del entorno
 ; \setsubject{csoc} ; \begin{document}...
% =====================================================================
\documentclass[aspectratio=169, 11pt]{beamer}
\usetheme{default}
\setbeamertemplate{navigation symbols}{}   % sin botones de navegación

\usepackage{fontspec}
\usepackage{polyglossia}\setdefaultlanguage{spanish}
\usepackage{tikz}
\usetikzlibrary{calc, positioning, arrows.meta, decorations.pathreplacing}
\usepackage{amsmath}

% --- TIPOGRAFÍA (idéntica familia que el libro) -----------------------
\setsansfont{Inter}[UprightFont=Inter-Regular, BoldFont=Inter-SemiBold,
  ItalicFont=Inter-Italic, Ligatures=TeX, Numbers=Lining]
\newfontfamily\Display{Inter Display}[Ligatures=TeX]
\newfontfamily\DisplayBlk{Inter Display Black}[Ligatures=TeX]
\newfontfamily\DisplaySB{Inter Display SemiBold}[Ligatures=TeX]
\newfontfamily\DisplayThin{Inter Display Thin}[Ligatures=TeX]

% --- COLOR (mismas 4 paletas que el libro) ----------------------------
\definecolor{mat} {RGB}{ 30,  64, 120}
\definecolor{cnat}{RGB}{ 16, 110,  80}
\definecolor{csoc}{RGB}{165,  75,  45}
\definecolor{eco} {RGB}{ 70,  55, 130}
\definecolor{ink}{RGB}{28,28,30}        % casi-negro para texto
\definecolor{paper}{RGB}{250,250,248}   % blanco cálido de fondo
\colorlet{subjectcolor}{csoc}
% \setsubject re-aplica los colores de beamer además de actualizar el alias
\newcommand{\setsubject}[1]{%
  \colorlet{subjectcolor}{#1}%
  \setbeamercolor{structure}{fg=#1}%
  \setbeamercolor{frametitle}{fg=#1,bg=}%
}

\setbeamercolor{normal text}{fg=ink, bg=paper}
\setbeamercolor{structure}{fg=subjectcolor}
\setbeamerfont{frametitle}{family=\DisplaySB}
\setbeamercolor{frametitle}{fg=subjectcolor, bg=}
\setbeamertemplate{frametitle}{%
  \vspace{0.6em}\hspace{0.2em}%
  {\DisplaySB\large\color{subjectcolor}\insertframetitle}\par
  \vspace{-0.2em}{\color{subjectcolor!30}\rule{\dimexpr\paperwidth-1.4em\relax}{1.2pt}}}

% Pie minimal: instituto · profe · tema · número en cuadro de color
\newcommand{\bmrtema}{}
\newcommand{\bmrprofe}{}   % ← \renewcommand{\bmrprofe}{Prof. Nombre} en cada beamer
\setbeamertemplate{footline}{%
  \hbox{}\vspace{2pt}\hspace{0.6em}%
  {\footnotesize\color{ink!45}ITSTZ%
   \ifx\bmrprofe\empty\else\space·\space\bmrprofe\fi}%
  \hfill{\footnotesize\color{ink!45}\bmrtema}\hspace{0.6em}%
  \tikz[baseline=-3pt]{\node[fill=subjectcolor, text=white, inner xsep=6pt,
     inner ysep=2.5pt, font=\DisplaySB\scriptsize]{\insertframenumber};}%
  \hspace{0.6em}\vspace{4pt}}

% =====================================================================
%  PORTADA y DIVISORES
% =====================================================================
\newcommand{\portada}[3]{% {kicker}{TÍTULO}{subtítulo}
{\setbeamertemplate{footline}{}
\begin{frame}[plain]
\begin{tikzpicture}[remember picture, overlay]
  \fill[paper] (current page.south west) rectangle (current page.north east);
  \node[anchor=north west, text=subjectcolor!12, font=\DisplayBlk]
    at ([xshift=-0.5cm,yshift=1.2cm]current page.north west)
    {\fontsize{190}{190}\selectfont ·};
  \node[anchor=west, text=subjectcolor, font=\DisplaySB\Large]
    at ([xshift=1.2cm, yshift=1.6cm]current page.west) {#1};
  \node[anchor=west, text=ink, font=\DisplayBlk, align=left]
    at ([xshift=1.15cm, yshift=0.2cm]current page.west)
    {\fontsize{40}{44}\selectfont #2};
  \fill[subjectcolor] ([xshift=1.2cm, yshift=-1.4cm]current page.west)
    rectangle ++(3.2cm, 0.16cm);
  \node[anchor=west, text=ink!60, font=\Display, align=left, text width=0.7\paperwidth]
    at ([xshift=1.2cm, yshift=-2.3cm]current page.west) {#3};
\end{tikzpicture}
\end{frame}}}

% Lámina divisoria: solo una frase grande centrada sobre color
\newcommand{\seccion}[1]{%
{\setbeamertemplate{footline}{}
\begin{frame}[plain]
\begin{tikzpicture}[remember picture, overlay]
  \fill[subjectcolor] (current page.south west) rectangle (current page.north east);
  \node[text=white, font=\DisplaySB, align=center, text width=0.85\paperwidth]
    at (current page.center) {\fontsize{30}{36}\selectfont #1};
\end{tikzpicture}
\end{frame}}}

% =====================================================================
%  ELEMENTOS DE CONTENIDO
% =====================================================================
% Una afirmación gigante que llena la lámina
\newcommand{\idea}[1]{%
{\setbeamertemplate{footline}{}
\begin{frame}[plain]
\begin{tikzpicture}[remember picture, overlay]
  \fill[paper] (current page.south west) rectangle (current page.north east);
  \fill[subjectcolor] (current page.north west)
    rectangle ([yshift=-0.32cm]current page.north east);
  \node[text=ink, font=\DisplaySB, align=center,
        text width=0.82\paperwidth]
    at (current page.center)
    {\fontsize{32}{44}\selectfont #1};
\end{tikzpicture}
\end{frame}}}

% Realce inline (fondo color suave, texto oscuro)
\newcommand{\resalta}[1]{\tikz[baseline=(R.base)]{\node[fill=subjectcolor!18,
  rounded corners=2pt, inner xsep=3pt, inner ysep=1pt](R){#1};}}
% Realce fuerte (fondo sólido, texto blanco)
\newcommand{\marca}[1]{\tikz[baseline=(R.base)]{\node[fill=subjectcolor,
  text=white, rounded corners=2pt, inner xsep=4pt, inner ysep=1.5pt](R){\textbf{#1}};}}

% Pregunta viva de apertura
\newcommand{\pregunta}[1]{%
\begin{frame}[plain]
\begin{tikzpicture}[remember picture,overlay]
  \fill[subjectcolor!7] (current page.south west) rectangle (current page.north east);
  \fill[subjectcolor] (current page.south west)
    rectangle ([xshift=0.45cm]current page.north west);
\end{tikzpicture}
\vfill
\hspace{1.5cm}\begin{minipage}{0.78\paperwidth}
  {\color{subjectcolor}\DisplayBlk\fontsize{54}{54}\selectfont ¿}\par
  \vspace{0.25em}
  {\Display\color{ink}\fontsize{22}{28}\selectfont #1}\par
  \vspace{0.4em}
  {\hfill\color{subjectcolor}\DisplayBlk\fontsize{38}{38}\selectfont ?}
\end{minipage}
\vfill
\end{frame}}

% Cierre / acción — label opcional (default: "DE LA QUEJA A LA ACCIÓN")
% Uso: \accion{texto}  o  \accion[ETIQUETA CUSTOM]{texto}
\newcommand{\accion}[2][DE LA QUEJA A LA ACCIÓN]{%
\begin{frame}[plain]
\begin{tikzpicture}[remember picture, overlay]
  \fill[subjectcolor!10] (current page.south west) rectangle (current page.north east);
  \node[anchor=west, text=subjectcolor, font=\DisplaySB]
    at ([xshift=1cm, yshift=2.2cm]current page.west)
    {\large #1};
  \node[anchor=west, text=ink, font=\Display, align=left, text width=0.78\paperwidth]
    at ([xshift=1cm, yshift=0cm]current page.west)
    {\fontsize{20}{26}\selectfont #2};
\end{tikzpicture}
\end{frame}}

% =====================================================================
%  LÍNEA DE TIEMPO HORIZONTAL
%  \begin{lineah}{año_ini}{año_fin}
%     \hito{posición 0-1}{etiqueta}   ...
%  \end{lineah}
% =====================================================================
\newenvironment{lineah}{%
  \begin{center}\begin{tikzpicture}[x=0.92\paperwidth, font=\footnotesize\sffamily]
  \draw[line width=2pt, subjectcolor, -{Stealth[length=5pt]}] (0,0) -- (1.02,0);
}{%
  \end{tikzpicture}\end{center}}
% hito: #1 = posición relativa 0..1 ; #2 = etiqueta (texto corto)
\newcommand{\hito}[2]{%
  \fill[subjectcolor] (#1,0) circle (3.2pt);
  \node[above, text=ink, align=center, text width=2.6cm, anchor=south,
        yshift=2pt] at (#1,0) {\bfseries #2};}
% hito bajo la línea (para alternar y no chocar)
\newcommand{\hitob}[2]{%
  \fill[subjectcolor] (#1,0) circle (3.2pt);
  \node[below, text=ink, align=center, text width=2.6cm, anchor=north,
        yshift=-2pt] at (#1,0) {#2};}

% =====================================================================
%  LÍNEA DE TIEMPO VERTICAL (ideal para listas cronológicas largas)
%  \begin{lineav} \hitov{año}{Título}{detalle} ... \end{lineav}
% =====================================================================
\newcounter{lvrow}
\newenvironment{lineav}{%
  \setcounter{lvrow}{0}%
  \begin{list}{}{\leftmargin=1.6cm \itemsep=0.55em \topsep=0.3em}%
  \renewcommand{\hitov}[3]{\stepcounter{lvrow}%
    \item[\tikz{\node[circle, fill=subjectcolor, text=white, inner sep=0pt,
       minimum size=1.5em, font=\DisplaySB\scriptsize]{\thelvrow};}]%
    {\DisplaySB\color{subjectcolor}##1}\quad{\bfseries ##2}\\[1pt]
    {\small\color{ink!75}##3}}%
}{%
  \end{list}}
\newcommand{\hitov}[3]{}  % placeholder; redefinido dentro del entorno

% =====================================================================
%  MAPA MENTAL — nodo central + ramas radiales
%  \begin{mapa}{NODO CENTRAL}
%     \rama{ángulo}{distancia cm}{texto de la rama}  ...
%  \end{mapa}
% =====================================================================
\newenvironment{mapa}[2][1]{%   % [escala opcional, def=1]{NODO CENTRAL}
  \begin{center}\begin{tikzpicture}[font=\sffamily\small, scale=#1, transform shape]
  \node[circle, fill=subjectcolor, text=white, font=\DisplaySB,
        align=center, text width=2.4cm, inner sep=4pt, minimum size=2.6cm]
        (CENTRO) at (0,0) {#2};
  \renewcommand{\rama}[3]{%
    \node[draw=subjectcolor!55, fill=subjectcolor!8, rounded corners=4pt,
          text=ink, align=center, text width=2.7cm, inner sep=4pt]
          (R##1) at (##1:##2) {##3};
    \draw[subjectcolor!55, line width=1.2pt] (CENTRO) -- (R##1);}%
}{%
  \end{tikzpicture}\end{center}}
\newcommand{\rama}[3]{}  % placeholder; redefinido dentro del entorno

%     % =====================================================================
%  TEMA · MATEMÁTICAS — variante del design system (BEAMER 16:9)
%  Identidad para decks de matemáticas (azul \setsubject{mat}):
%    · FIX del bug de preflight §5/§6: las láminas a sangre del design
%      system (\portada \seccion \idea \pregunta \accion) usan
%      remember picture + current page y salen CORRIDAS hacia arriba con
%      compile_quiet.sh. Aquí se redefinen con `background canvas` +
%      coordenadas absolutas (misma técnica probada del tema eco/cnat).
%    · carga lmodern/colortbl/enumitem (que design/beamer.tex no trae) para
%      math escalable en titulares, \rowcolor en tablas y itemize[clave=valor].
%    · \DisplaySB con BoldFont explícito → \textbf dentro de \idea no avisa.
%    · portada con el personaje "Dr. Cuack" matemático (opcional; se dibuja
%      solo si existe el PNG).
%
%  INVOCACIÓN (tras % =====================================================================
%  DESIGN SYSTEM BEAMER v5 — ITSTZ · 16:9
%  Coherente con el libro: Inter Display (títulos) + Inter (cuerpo),
%  un color de asignatura, fondo claro, texto GRANDE y contrastado.
%  Filosofía slide: una idea por lámina, mucho aire, cero relleno.
%
%  Vocabulario:
%    \portada{kicker}{TÍTULO}{subtítulo}
%    \seccion{texto}                  lámina divisoria a pantalla completa
%    \idea{TEXTO GRANDE}              una afirmación que llena la slide
%    \resalta{texto}                  realce inline (fondo color)
%    frames de contenido normales con \frametitle
%    \begin{lineah} \hito{x}{etiqueta} ... \end{lineah}   timeline horizontal
%    \begin{lineav} \hitov{x}{título}{texto} ... \end{lineav} timeline vertical
%    \begin{mapa}{NODO CENTRAL} \rama{ang}{dist}{texto} ... \end{mapa}
%    \pregunta{texto}                 pregunta viva (apertura)
%    \accion{texto}                   cierre "de la queja a la acción"
%  Uso: % =====================================================================
%  DESIGN SYSTEM BEAMER v5 — ITSTZ · 16:9
%  Coherente con el libro: Inter Display (títulos) + Inter (cuerpo),
%  un color de asignatura, fondo claro, texto GRANDE y contrastado.
%  Filosofía slide: una idea por lámina, mucho aire, cero relleno.
%
%  Vocabulario:
%    \portada{kicker}{TÍTULO}{subtítulo}
%    \seccion{texto}                  lámina divisoria a pantalla completa
%    \idea{TEXTO GRANDE}              una afirmación que llena la slide
%    \resalta{texto}                  realce inline (fondo color)
%    frames de contenido normales con \frametitle
%    \begin{lineah} \hito{x}{etiqueta} ... \end{lineah}   timeline horizontal
%    \begin{lineav} \hitov{x}{título}{texto} ... \end{lineav} timeline vertical
%    \begin{mapa}{NODO CENTRAL} \rama{ang}{dist}{texto} ... \end{mapa}
%    \pregunta{texto}                 pregunta viva (apertura)
%    \accion{texto}                   cierre "de la queja a la acción"
%  Uso: \input{design/beamer.tex} ; \setsubject{csoc} ; \begin{document}...
% =====================================================================
\documentclass[aspectratio=169, 11pt]{beamer}
\usetheme{default}
\setbeamertemplate{navigation symbols}{}   % sin botones de navegación

\usepackage{fontspec}
\usepackage{polyglossia}\setdefaultlanguage{spanish}
\usepackage{tikz}
\usetikzlibrary{calc, positioning, arrows.meta, decorations.pathreplacing}
\usepackage{amsmath}

% --- TIPOGRAFÍA (idéntica familia que el libro) -----------------------
\setsansfont{Inter}[UprightFont=Inter-Regular, BoldFont=Inter-SemiBold,
  ItalicFont=Inter-Italic, Ligatures=TeX, Numbers=Lining]
\newfontfamily\Display{Inter Display}[Ligatures=TeX]
\newfontfamily\DisplayBlk{Inter Display Black}[Ligatures=TeX]
\newfontfamily\DisplaySB{Inter Display SemiBold}[Ligatures=TeX]
\newfontfamily\DisplayThin{Inter Display Thin}[Ligatures=TeX]

% --- COLOR (mismas 4 paletas que el libro) ----------------------------
\definecolor{mat} {RGB}{ 30,  64, 120}
\definecolor{cnat}{RGB}{ 16, 110,  80}
\definecolor{csoc}{RGB}{165,  75,  45}
\definecolor{eco} {RGB}{ 70,  55, 130}
\definecolor{ink}{RGB}{28,28,30}        % casi-negro para texto
\definecolor{paper}{RGB}{250,250,248}   % blanco cálido de fondo
\colorlet{subjectcolor}{csoc}
% \setsubject re-aplica los colores de beamer además de actualizar el alias
\newcommand{\setsubject}[1]{%
  \colorlet{subjectcolor}{#1}%
  \setbeamercolor{structure}{fg=#1}%
  \setbeamercolor{frametitle}{fg=#1,bg=}%
}

\setbeamercolor{normal text}{fg=ink, bg=paper}
\setbeamercolor{structure}{fg=subjectcolor}
\setbeamerfont{frametitle}{family=\DisplaySB}
\setbeamercolor{frametitle}{fg=subjectcolor, bg=}
\setbeamertemplate{frametitle}{%
  \vspace{0.6em}\hspace{0.2em}%
  {\DisplaySB\large\color{subjectcolor}\insertframetitle}\par
  \vspace{-0.2em}{\color{subjectcolor!30}\rule{\dimexpr\paperwidth-1.4em\relax}{1.2pt}}}

% Pie minimal: instituto · profe · tema · número en cuadro de color
\newcommand{\bmrtema}{}
\newcommand{\bmrprofe}{}   % ← \renewcommand{\bmrprofe}{Prof. Nombre} en cada beamer
\setbeamertemplate{footline}{%
  \hbox{}\vspace{2pt}\hspace{0.6em}%
  {\footnotesize\color{ink!45}ITSTZ%
   \ifx\bmrprofe\empty\else\space·\space\bmrprofe\fi}%
  \hfill{\footnotesize\color{ink!45}\bmrtema}\hspace{0.6em}%
  \tikz[baseline=-3pt]{\node[fill=subjectcolor, text=white, inner xsep=6pt,
     inner ysep=2.5pt, font=\DisplaySB\scriptsize]{\insertframenumber};}%
  \hspace{0.6em}\vspace{4pt}}

% =====================================================================
%  PORTADA y DIVISORES
% =====================================================================
\newcommand{\portada}[3]{% {kicker}{TÍTULO}{subtítulo}
{\setbeamertemplate{footline}{}
\begin{frame}[plain]
\begin{tikzpicture}[remember picture, overlay]
  \fill[paper] (current page.south west) rectangle (current page.north east);
  \node[anchor=north west, text=subjectcolor!12, font=\DisplayBlk]
    at ([xshift=-0.5cm,yshift=1.2cm]current page.north west)
    {\fontsize{190}{190}\selectfont ·};
  \node[anchor=west, text=subjectcolor, font=\DisplaySB\Large]
    at ([xshift=1.2cm, yshift=1.6cm]current page.west) {#1};
  \node[anchor=west, text=ink, font=\DisplayBlk, align=left]
    at ([xshift=1.15cm, yshift=0.2cm]current page.west)
    {\fontsize{40}{44}\selectfont #2};
  \fill[subjectcolor] ([xshift=1.2cm, yshift=-1.4cm]current page.west)
    rectangle ++(3.2cm, 0.16cm);
  \node[anchor=west, text=ink!60, font=\Display, align=left, text width=0.7\paperwidth]
    at ([xshift=1.2cm, yshift=-2.3cm]current page.west) {#3};
\end{tikzpicture}
\end{frame}}}

% Lámina divisoria: solo una frase grande centrada sobre color
\newcommand{\seccion}[1]{%
{\setbeamertemplate{footline}{}
\begin{frame}[plain]
\begin{tikzpicture}[remember picture, overlay]
  \fill[subjectcolor] (current page.south west) rectangle (current page.north east);
  \node[text=white, font=\DisplaySB, align=center, text width=0.85\paperwidth]
    at (current page.center) {\fontsize{30}{36}\selectfont #1};
\end{tikzpicture}
\end{frame}}}

% =====================================================================
%  ELEMENTOS DE CONTENIDO
% =====================================================================
% Una afirmación gigante que llena la lámina
\newcommand{\idea}[1]{%
{\setbeamertemplate{footline}{}
\begin{frame}[plain]
\begin{tikzpicture}[remember picture, overlay]
  \fill[paper] (current page.south west) rectangle (current page.north east);
  \fill[subjectcolor] (current page.north west)
    rectangle ([yshift=-0.32cm]current page.north east);
  \node[text=ink, font=\DisplaySB, align=center,
        text width=0.82\paperwidth]
    at (current page.center)
    {\fontsize{32}{44}\selectfont #1};
\end{tikzpicture}
\end{frame}}}

% Realce inline (fondo color suave, texto oscuro)
\newcommand{\resalta}[1]{\tikz[baseline=(R.base)]{\node[fill=subjectcolor!18,
  rounded corners=2pt, inner xsep=3pt, inner ysep=1pt](R){#1};}}
% Realce fuerte (fondo sólido, texto blanco)
\newcommand{\marca}[1]{\tikz[baseline=(R.base)]{\node[fill=subjectcolor,
  text=white, rounded corners=2pt, inner xsep=4pt, inner ysep=1.5pt](R){\textbf{#1}};}}

% Pregunta viva de apertura
\newcommand{\pregunta}[1]{%
\begin{frame}[plain]
\begin{tikzpicture}[remember picture,overlay]
  \fill[subjectcolor!7] (current page.south west) rectangle (current page.north east);
  \fill[subjectcolor] (current page.south west)
    rectangle ([xshift=0.45cm]current page.north west);
\end{tikzpicture}
\vfill
\hspace{1.5cm}\begin{minipage}{0.78\paperwidth}
  {\color{subjectcolor}\DisplayBlk\fontsize{54}{54}\selectfont ¿}\par
  \vspace{0.25em}
  {\Display\color{ink}\fontsize{22}{28}\selectfont #1}\par
  \vspace{0.4em}
  {\hfill\color{subjectcolor}\DisplayBlk\fontsize{38}{38}\selectfont ?}
\end{minipage}
\vfill
\end{frame}}

% Cierre / acción — label opcional (default: "DE LA QUEJA A LA ACCIÓN")
% Uso: \accion{texto}  o  \accion[ETIQUETA CUSTOM]{texto}
\newcommand{\accion}[2][DE LA QUEJA A LA ACCIÓN]{%
\begin{frame}[plain]
\begin{tikzpicture}[remember picture, overlay]
  \fill[subjectcolor!10] (current page.south west) rectangle (current page.north east);
  \node[anchor=west, text=subjectcolor, font=\DisplaySB]
    at ([xshift=1cm, yshift=2.2cm]current page.west)
    {\large #1};
  \node[anchor=west, text=ink, font=\Display, align=left, text width=0.78\paperwidth]
    at ([xshift=1cm, yshift=0cm]current page.west)
    {\fontsize{20}{26}\selectfont #2};
\end{tikzpicture}
\end{frame}}

% =====================================================================
%  LÍNEA DE TIEMPO HORIZONTAL
%  \begin{lineah}{año_ini}{año_fin}
%     \hito{posición 0-1}{etiqueta}   ...
%  \end{lineah}
% =====================================================================
\newenvironment{lineah}{%
  \begin{center}\begin{tikzpicture}[x=0.92\paperwidth, font=\footnotesize\sffamily]
  \draw[line width=2pt, subjectcolor, -{Stealth[length=5pt]}] (0,0) -- (1.02,0);
}{%
  \end{tikzpicture}\end{center}}
% hito: #1 = posición relativa 0..1 ; #2 = etiqueta (texto corto)
\newcommand{\hito}[2]{%
  \fill[subjectcolor] (#1,0) circle (3.2pt);
  \node[above, text=ink, align=center, text width=2.6cm, anchor=south,
        yshift=2pt] at (#1,0) {\bfseries #2};}
% hito bajo la línea (para alternar y no chocar)
\newcommand{\hitob}[2]{%
  \fill[subjectcolor] (#1,0) circle (3.2pt);
  \node[below, text=ink, align=center, text width=2.6cm, anchor=north,
        yshift=-2pt] at (#1,0) {#2};}

% =====================================================================
%  LÍNEA DE TIEMPO VERTICAL (ideal para listas cronológicas largas)
%  \begin{lineav} \hitov{año}{Título}{detalle} ... \end{lineav}
% =====================================================================
\newcounter{lvrow}
\newenvironment{lineav}{%
  \setcounter{lvrow}{0}%
  \begin{list}{}{\leftmargin=1.6cm \itemsep=0.55em \topsep=0.3em}%
  \renewcommand{\hitov}[3]{\stepcounter{lvrow}%
    \item[\tikz{\node[circle, fill=subjectcolor, text=white, inner sep=0pt,
       minimum size=1.5em, font=\DisplaySB\scriptsize]{\thelvrow};}]%
    {\DisplaySB\color{subjectcolor}##1}\quad{\bfseries ##2}\\[1pt]
    {\small\color{ink!75}##3}}%
}{%
  \end{list}}
\newcommand{\hitov}[3]{}  % placeholder; redefinido dentro del entorno

% =====================================================================
%  MAPA MENTAL — nodo central + ramas radiales
%  \begin{mapa}{NODO CENTRAL}
%     \rama{ángulo}{distancia cm}{texto de la rama}  ...
%  \end{mapa}
% =====================================================================
\newenvironment{mapa}[2][1]{%   % [escala opcional, def=1]{NODO CENTRAL}
  \begin{center}\begin{tikzpicture}[font=\sffamily\small, scale=#1, transform shape]
  \node[circle, fill=subjectcolor, text=white, font=\DisplaySB,
        align=center, text width=2.4cm, inner sep=4pt, minimum size=2.6cm]
        (CENTRO) at (0,0) {#2};
  \renewcommand{\rama}[3]{%
    \node[draw=subjectcolor!55, fill=subjectcolor!8, rounded corners=4pt,
          text=ink, align=center, text width=2.7cm, inner sep=4pt]
          (R##1) at (##1:##2) {##3};
    \draw[subjectcolor!55, line width=1.2pt] (CENTRO) -- (R##1);}%
}{%
  \end{tikzpicture}\end{center}}
\newcommand{\rama}[3]{}  % placeholder; redefinido dentro del entorno
 ; \setsubject{csoc} ; \begin{document}...
% =====================================================================
\documentclass[aspectratio=169, 11pt]{beamer}
\usetheme{default}
\setbeamertemplate{navigation symbols}{}   % sin botones de navegación

\usepackage{fontspec}
\usepackage{polyglossia}\setdefaultlanguage{spanish}
\usepackage{tikz}
\usetikzlibrary{calc, positioning, arrows.meta, decorations.pathreplacing}
\usepackage{amsmath}

% --- TIPOGRAFÍA (idéntica familia que el libro) -----------------------
\setsansfont{Inter}[UprightFont=Inter-Regular, BoldFont=Inter-SemiBold,
  ItalicFont=Inter-Italic, Ligatures=TeX, Numbers=Lining]
\newfontfamily\Display{Inter Display}[Ligatures=TeX]
\newfontfamily\DisplayBlk{Inter Display Black}[Ligatures=TeX]
\newfontfamily\DisplaySB{Inter Display SemiBold}[Ligatures=TeX]
\newfontfamily\DisplayThin{Inter Display Thin}[Ligatures=TeX]

% --- COLOR (mismas 4 paletas que el libro) ----------------------------
\definecolor{mat} {RGB}{ 30,  64, 120}
\definecolor{cnat}{RGB}{ 16, 110,  80}
\definecolor{csoc}{RGB}{165,  75,  45}
\definecolor{eco} {RGB}{ 70,  55, 130}
\definecolor{ink}{RGB}{28,28,30}        % casi-negro para texto
\definecolor{paper}{RGB}{250,250,248}   % blanco cálido de fondo
\colorlet{subjectcolor}{csoc}
% \setsubject re-aplica los colores de beamer además de actualizar el alias
\newcommand{\setsubject}[1]{%
  \colorlet{subjectcolor}{#1}%
  \setbeamercolor{structure}{fg=#1}%
  \setbeamercolor{frametitle}{fg=#1,bg=}%
}

\setbeamercolor{normal text}{fg=ink, bg=paper}
\setbeamercolor{structure}{fg=subjectcolor}
\setbeamerfont{frametitle}{family=\DisplaySB}
\setbeamercolor{frametitle}{fg=subjectcolor, bg=}
\setbeamertemplate{frametitle}{%
  \vspace{0.6em}\hspace{0.2em}%
  {\DisplaySB\large\color{subjectcolor}\insertframetitle}\par
  \vspace{-0.2em}{\color{subjectcolor!30}\rule{\dimexpr\paperwidth-1.4em\relax}{1.2pt}}}

% Pie minimal: instituto · profe · tema · número en cuadro de color
\newcommand{\bmrtema}{}
\newcommand{\bmrprofe}{}   % ← \renewcommand{\bmrprofe}{Prof. Nombre} en cada beamer
\setbeamertemplate{footline}{%
  \hbox{}\vspace{2pt}\hspace{0.6em}%
  {\footnotesize\color{ink!45}ITSTZ%
   \ifx\bmrprofe\empty\else\space·\space\bmrprofe\fi}%
  \hfill{\footnotesize\color{ink!45}\bmrtema}\hspace{0.6em}%
  \tikz[baseline=-3pt]{\node[fill=subjectcolor, text=white, inner xsep=6pt,
     inner ysep=2.5pt, font=\DisplaySB\scriptsize]{\insertframenumber};}%
  \hspace{0.6em}\vspace{4pt}}

% =====================================================================
%  PORTADA y DIVISORES
% =====================================================================
\newcommand{\portada}[3]{% {kicker}{TÍTULO}{subtítulo}
{\setbeamertemplate{footline}{}
\begin{frame}[plain]
\begin{tikzpicture}[remember picture, overlay]
  \fill[paper] (current page.south west) rectangle (current page.north east);
  \node[anchor=north west, text=subjectcolor!12, font=\DisplayBlk]
    at ([xshift=-0.5cm,yshift=1.2cm]current page.north west)
    {\fontsize{190}{190}\selectfont ·};
  \node[anchor=west, text=subjectcolor, font=\DisplaySB\Large]
    at ([xshift=1.2cm, yshift=1.6cm]current page.west) {#1};
  \node[anchor=west, text=ink, font=\DisplayBlk, align=left]
    at ([xshift=1.15cm, yshift=0.2cm]current page.west)
    {\fontsize{40}{44}\selectfont #2};
  \fill[subjectcolor] ([xshift=1.2cm, yshift=-1.4cm]current page.west)
    rectangle ++(3.2cm, 0.16cm);
  \node[anchor=west, text=ink!60, font=\Display, align=left, text width=0.7\paperwidth]
    at ([xshift=1.2cm, yshift=-2.3cm]current page.west) {#3};
\end{tikzpicture}
\end{frame}}}

% Lámina divisoria: solo una frase grande centrada sobre color
\newcommand{\seccion}[1]{%
{\setbeamertemplate{footline}{}
\begin{frame}[plain]
\begin{tikzpicture}[remember picture, overlay]
  \fill[subjectcolor] (current page.south west) rectangle (current page.north east);
  \node[text=white, font=\DisplaySB, align=center, text width=0.85\paperwidth]
    at (current page.center) {\fontsize{30}{36}\selectfont #1};
\end{tikzpicture}
\end{frame}}}

% =====================================================================
%  ELEMENTOS DE CONTENIDO
% =====================================================================
% Una afirmación gigante que llena la lámina
\newcommand{\idea}[1]{%
{\setbeamertemplate{footline}{}
\begin{frame}[plain]
\begin{tikzpicture}[remember picture, overlay]
  \fill[paper] (current page.south west) rectangle (current page.north east);
  \fill[subjectcolor] (current page.north west)
    rectangle ([yshift=-0.32cm]current page.north east);
  \node[text=ink, font=\DisplaySB, align=center,
        text width=0.82\paperwidth]
    at (current page.center)
    {\fontsize{32}{44}\selectfont #1};
\end{tikzpicture}
\end{frame}}}

% Realce inline (fondo color suave, texto oscuro)
\newcommand{\resalta}[1]{\tikz[baseline=(R.base)]{\node[fill=subjectcolor!18,
  rounded corners=2pt, inner xsep=3pt, inner ysep=1pt](R){#1};}}
% Realce fuerte (fondo sólido, texto blanco)
\newcommand{\marca}[1]{\tikz[baseline=(R.base)]{\node[fill=subjectcolor,
  text=white, rounded corners=2pt, inner xsep=4pt, inner ysep=1.5pt](R){\textbf{#1}};}}

% Pregunta viva de apertura
\newcommand{\pregunta}[1]{%
\begin{frame}[plain]
\begin{tikzpicture}[remember picture,overlay]
  \fill[subjectcolor!7] (current page.south west) rectangle (current page.north east);
  \fill[subjectcolor] (current page.south west)
    rectangle ([xshift=0.45cm]current page.north west);
\end{tikzpicture}
\vfill
\hspace{1.5cm}\begin{minipage}{0.78\paperwidth}
  {\color{subjectcolor}\DisplayBlk\fontsize{54}{54}\selectfont ¿}\par
  \vspace{0.25em}
  {\Display\color{ink}\fontsize{22}{28}\selectfont #1}\par
  \vspace{0.4em}
  {\hfill\color{subjectcolor}\DisplayBlk\fontsize{38}{38}\selectfont ?}
\end{minipage}
\vfill
\end{frame}}

% Cierre / acción — label opcional (default: "DE LA QUEJA A LA ACCIÓN")
% Uso: \accion{texto}  o  \accion[ETIQUETA CUSTOM]{texto}
\newcommand{\accion}[2][DE LA QUEJA A LA ACCIÓN]{%
\begin{frame}[plain]
\begin{tikzpicture}[remember picture, overlay]
  \fill[subjectcolor!10] (current page.south west) rectangle (current page.north east);
  \node[anchor=west, text=subjectcolor, font=\DisplaySB]
    at ([xshift=1cm, yshift=2.2cm]current page.west)
    {\large #1};
  \node[anchor=west, text=ink, font=\Display, align=left, text width=0.78\paperwidth]
    at ([xshift=1cm, yshift=0cm]current page.west)
    {\fontsize{20}{26}\selectfont #2};
\end{tikzpicture}
\end{frame}}

% =====================================================================
%  LÍNEA DE TIEMPO HORIZONTAL
%  \begin{lineah}{año_ini}{año_fin}
%     \hito{posición 0-1}{etiqueta}   ...
%  \end{lineah}
% =====================================================================
\newenvironment{lineah}{%
  \begin{center}\begin{tikzpicture}[x=0.92\paperwidth, font=\footnotesize\sffamily]
  \draw[line width=2pt, subjectcolor, -{Stealth[length=5pt]}] (0,0) -- (1.02,0);
}{%
  \end{tikzpicture}\end{center}}
% hito: #1 = posición relativa 0..1 ; #2 = etiqueta (texto corto)
\newcommand{\hito}[2]{%
  \fill[subjectcolor] (#1,0) circle (3.2pt);
  \node[above, text=ink, align=center, text width=2.6cm, anchor=south,
        yshift=2pt] at (#1,0) {\bfseries #2};}
% hito bajo la línea (para alternar y no chocar)
\newcommand{\hitob}[2]{%
  \fill[subjectcolor] (#1,0) circle (3.2pt);
  \node[below, text=ink, align=center, text width=2.6cm, anchor=north,
        yshift=-2pt] at (#1,0) {#2};}

% =====================================================================
%  LÍNEA DE TIEMPO VERTICAL (ideal para listas cronológicas largas)
%  \begin{lineav} \hitov{año}{Título}{detalle} ... \end{lineav}
% =====================================================================
\newcounter{lvrow}
\newenvironment{lineav}{%
  \setcounter{lvrow}{0}%
  \begin{list}{}{\leftmargin=1.6cm \itemsep=0.55em \topsep=0.3em}%
  \renewcommand{\hitov}[3]{\stepcounter{lvrow}%
    \item[\tikz{\node[circle, fill=subjectcolor, text=white, inner sep=0pt,
       minimum size=1.5em, font=\DisplaySB\scriptsize]{\thelvrow};}]%
    {\DisplaySB\color{subjectcolor}##1}\quad{\bfseries ##2}\\[1pt]
    {\small\color{ink!75}##3}}%
}{%
  \end{list}}
\newcommand{\hitov}[3]{}  % placeholder; redefinido dentro del entorno

% =====================================================================
%  MAPA MENTAL — nodo central + ramas radiales
%  \begin{mapa}{NODO CENTRAL}
%     \rama{ángulo}{distancia cm}{texto de la rama}  ...
%  \end{mapa}
% =====================================================================
\newenvironment{mapa}[2][1]{%   % [escala opcional, def=1]{NODO CENTRAL}
  \begin{center}\begin{tikzpicture}[font=\sffamily\small, scale=#1, transform shape]
  \node[circle, fill=subjectcolor, text=white, font=\DisplaySB,
        align=center, text width=2.4cm, inner sep=4pt, minimum size=2.6cm]
        (CENTRO) at (0,0) {#2};
  \renewcommand{\rama}[3]{%
    \node[draw=subjectcolor!55, fill=subjectcolor!8, rounded corners=4pt,
          text=ink, align=center, text width=2.7cm, inner sep=4pt]
          (R##1) at (##1:##2) {##3};
    \draw[subjectcolor!55, line width=1.2pt] (CENTRO) -- (R##1);}%
}{%
  \end{tikzpicture}\end{center}}
\newcommand{\rama}[3]{}  % placeholder; redefinido dentro del entorno
):
%     % =====================================================================
%  DESIGN SYSTEM BEAMER v5 — ITSTZ · 16:9
%  Coherente con el libro: Inter Display (títulos) + Inter (cuerpo),
%  un color de asignatura, fondo claro, texto GRANDE y contrastado.
%  Filosofía slide: una idea por lámina, mucho aire, cero relleno.
%
%  Vocabulario:
%    \portada{kicker}{TÍTULO}{subtítulo}
%    \seccion{texto}                  lámina divisoria a pantalla completa
%    \idea{TEXTO GRANDE}              una afirmación que llena la slide
%    \resalta{texto}                  realce inline (fondo color)
%    frames de contenido normales con \frametitle
%    \begin{lineah} \hito{x}{etiqueta} ... \end{lineah}   timeline horizontal
%    \begin{lineav} \hitov{x}{título}{texto} ... \end{lineav} timeline vertical
%    \begin{mapa}{NODO CENTRAL} \rama{ang}{dist}{texto} ... \end{mapa}
%    \pregunta{texto}                 pregunta viva (apertura)
%    \accion{texto}                   cierre "de la queja a la acción"
%  Uso: % =====================================================================
%  DESIGN SYSTEM BEAMER v5 — ITSTZ · 16:9
%  Coherente con el libro: Inter Display (títulos) + Inter (cuerpo),
%  un color de asignatura, fondo claro, texto GRANDE y contrastado.
%  Filosofía slide: una idea por lámina, mucho aire, cero relleno.
%
%  Vocabulario:
%    \portada{kicker}{TÍTULO}{subtítulo}
%    \seccion{texto}                  lámina divisoria a pantalla completa
%    \idea{TEXTO GRANDE}              una afirmación que llena la slide
%    \resalta{texto}                  realce inline (fondo color)
%    frames de contenido normales con \frametitle
%    \begin{lineah} \hito{x}{etiqueta} ... \end{lineah}   timeline horizontal
%    \begin{lineav} \hitov{x}{título}{texto} ... \end{lineav} timeline vertical
%    \begin{mapa}{NODO CENTRAL} \rama{ang}{dist}{texto} ... \end{mapa}
%    \pregunta{texto}                 pregunta viva (apertura)
%    \accion{texto}                   cierre "de la queja a la acción"
%  Uso: \input{design/beamer.tex} ; \setsubject{csoc} ; \begin{document}...
% =====================================================================
\documentclass[aspectratio=169, 11pt]{beamer}
\usetheme{default}
\setbeamertemplate{navigation symbols}{}   % sin botones de navegación

\usepackage{fontspec}
\usepackage{polyglossia}\setdefaultlanguage{spanish}
\usepackage{tikz}
\usetikzlibrary{calc, positioning, arrows.meta, decorations.pathreplacing}
\usepackage{amsmath}

% --- TIPOGRAFÍA (idéntica familia que el libro) -----------------------
\setsansfont{Inter}[UprightFont=Inter-Regular, BoldFont=Inter-SemiBold,
  ItalicFont=Inter-Italic, Ligatures=TeX, Numbers=Lining]
\newfontfamily\Display{Inter Display}[Ligatures=TeX]
\newfontfamily\DisplayBlk{Inter Display Black}[Ligatures=TeX]
\newfontfamily\DisplaySB{Inter Display SemiBold}[Ligatures=TeX]
\newfontfamily\DisplayThin{Inter Display Thin}[Ligatures=TeX]

% --- COLOR (mismas 4 paletas que el libro) ----------------------------
\definecolor{mat} {RGB}{ 30,  64, 120}
\definecolor{cnat}{RGB}{ 16, 110,  80}
\definecolor{csoc}{RGB}{165,  75,  45}
\definecolor{eco} {RGB}{ 70,  55, 130}
\definecolor{ink}{RGB}{28,28,30}        % casi-negro para texto
\definecolor{paper}{RGB}{250,250,248}   % blanco cálido de fondo
\colorlet{subjectcolor}{csoc}
% \setsubject re-aplica los colores de beamer además de actualizar el alias
\newcommand{\setsubject}[1]{%
  \colorlet{subjectcolor}{#1}%
  \setbeamercolor{structure}{fg=#1}%
  \setbeamercolor{frametitle}{fg=#1,bg=}%
}

\setbeamercolor{normal text}{fg=ink, bg=paper}
\setbeamercolor{structure}{fg=subjectcolor}
\setbeamerfont{frametitle}{family=\DisplaySB}
\setbeamercolor{frametitle}{fg=subjectcolor, bg=}
\setbeamertemplate{frametitle}{%
  \vspace{0.6em}\hspace{0.2em}%
  {\DisplaySB\large\color{subjectcolor}\insertframetitle}\par
  \vspace{-0.2em}{\color{subjectcolor!30}\rule{\dimexpr\paperwidth-1.4em\relax}{1.2pt}}}

% Pie minimal: instituto · profe · tema · número en cuadro de color
\newcommand{\bmrtema}{}
\newcommand{\bmrprofe}{}   % ← \renewcommand{\bmrprofe}{Prof. Nombre} en cada beamer
\setbeamertemplate{footline}{%
  \hbox{}\vspace{2pt}\hspace{0.6em}%
  {\footnotesize\color{ink!45}ITSTZ%
   \ifx\bmrprofe\empty\else\space·\space\bmrprofe\fi}%
  \hfill{\footnotesize\color{ink!45}\bmrtema}\hspace{0.6em}%
  \tikz[baseline=-3pt]{\node[fill=subjectcolor, text=white, inner xsep=6pt,
     inner ysep=2.5pt, font=\DisplaySB\scriptsize]{\insertframenumber};}%
  \hspace{0.6em}\vspace{4pt}}

% =====================================================================
%  PORTADA y DIVISORES
% =====================================================================
\newcommand{\portada}[3]{% {kicker}{TÍTULO}{subtítulo}
{\setbeamertemplate{footline}{}
\begin{frame}[plain]
\begin{tikzpicture}[remember picture, overlay]
  \fill[paper] (current page.south west) rectangle (current page.north east);
  \node[anchor=north west, text=subjectcolor!12, font=\DisplayBlk]
    at ([xshift=-0.5cm,yshift=1.2cm]current page.north west)
    {\fontsize{190}{190}\selectfont ·};
  \node[anchor=west, text=subjectcolor, font=\DisplaySB\Large]
    at ([xshift=1.2cm, yshift=1.6cm]current page.west) {#1};
  \node[anchor=west, text=ink, font=\DisplayBlk, align=left]
    at ([xshift=1.15cm, yshift=0.2cm]current page.west)
    {\fontsize{40}{44}\selectfont #2};
  \fill[subjectcolor] ([xshift=1.2cm, yshift=-1.4cm]current page.west)
    rectangle ++(3.2cm, 0.16cm);
  \node[anchor=west, text=ink!60, font=\Display, align=left, text width=0.7\paperwidth]
    at ([xshift=1.2cm, yshift=-2.3cm]current page.west) {#3};
\end{tikzpicture}
\end{frame}}}

% Lámina divisoria: solo una frase grande centrada sobre color
\newcommand{\seccion}[1]{%
{\setbeamertemplate{footline}{}
\begin{frame}[plain]
\begin{tikzpicture}[remember picture, overlay]
  \fill[subjectcolor] (current page.south west) rectangle (current page.north east);
  \node[text=white, font=\DisplaySB, align=center, text width=0.85\paperwidth]
    at (current page.center) {\fontsize{30}{36}\selectfont #1};
\end{tikzpicture}
\end{frame}}}

% =====================================================================
%  ELEMENTOS DE CONTENIDO
% =====================================================================
% Una afirmación gigante que llena la lámina
\newcommand{\idea}[1]{%
{\setbeamertemplate{footline}{}
\begin{frame}[plain]
\begin{tikzpicture}[remember picture, overlay]
  \fill[paper] (current page.south west) rectangle (current page.north east);
  \fill[subjectcolor] (current page.north west)
    rectangle ([yshift=-0.32cm]current page.north east);
  \node[text=ink, font=\DisplaySB, align=center,
        text width=0.82\paperwidth]
    at (current page.center)
    {\fontsize{32}{44}\selectfont #1};
\end{tikzpicture}
\end{frame}}}

% Realce inline (fondo color suave, texto oscuro)
\newcommand{\resalta}[1]{\tikz[baseline=(R.base)]{\node[fill=subjectcolor!18,
  rounded corners=2pt, inner xsep=3pt, inner ysep=1pt](R){#1};}}
% Realce fuerte (fondo sólido, texto blanco)
\newcommand{\marca}[1]{\tikz[baseline=(R.base)]{\node[fill=subjectcolor,
  text=white, rounded corners=2pt, inner xsep=4pt, inner ysep=1.5pt](R){\textbf{#1}};}}

% Pregunta viva de apertura
\newcommand{\pregunta}[1]{%
\begin{frame}[plain]
\begin{tikzpicture}[remember picture,overlay]
  \fill[subjectcolor!7] (current page.south west) rectangle (current page.north east);
  \fill[subjectcolor] (current page.south west)
    rectangle ([xshift=0.45cm]current page.north west);
\end{tikzpicture}
\vfill
\hspace{1.5cm}\begin{minipage}{0.78\paperwidth}
  {\color{subjectcolor}\DisplayBlk\fontsize{54}{54}\selectfont ¿}\par
  \vspace{0.25em}
  {\Display\color{ink}\fontsize{22}{28}\selectfont #1}\par
  \vspace{0.4em}
  {\hfill\color{subjectcolor}\DisplayBlk\fontsize{38}{38}\selectfont ?}
\end{minipage}
\vfill
\end{frame}}

% Cierre / acción — label opcional (default: "DE LA QUEJA A LA ACCIÓN")
% Uso: \accion{texto}  o  \accion[ETIQUETA CUSTOM]{texto}
\newcommand{\accion}[2][DE LA QUEJA A LA ACCIÓN]{%
\begin{frame}[plain]
\begin{tikzpicture}[remember picture, overlay]
  \fill[subjectcolor!10] (current page.south west) rectangle (current page.north east);
  \node[anchor=west, text=subjectcolor, font=\DisplaySB]
    at ([xshift=1cm, yshift=2.2cm]current page.west)
    {\large #1};
  \node[anchor=west, text=ink, font=\Display, align=left, text width=0.78\paperwidth]
    at ([xshift=1cm, yshift=0cm]current page.west)
    {\fontsize{20}{26}\selectfont #2};
\end{tikzpicture}
\end{frame}}

% =====================================================================
%  LÍNEA DE TIEMPO HORIZONTAL
%  \begin{lineah}{año_ini}{año_fin}
%     \hito{posición 0-1}{etiqueta}   ...
%  \end{lineah}
% =====================================================================
\newenvironment{lineah}{%
  \begin{center}\begin{tikzpicture}[x=0.92\paperwidth, font=\footnotesize\sffamily]
  \draw[line width=2pt, subjectcolor, -{Stealth[length=5pt]}] (0,0) -- (1.02,0);
}{%
  \end{tikzpicture}\end{center}}
% hito: #1 = posición relativa 0..1 ; #2 = etiqueta (texto corto)
\newcommand{\hito}[2]{%
  \fill[subjectcolor] (#1,0) circle (3.2pt);
  \node[above, text=ink, align=center, text width=2.6cm, anchor=south,
        yshift=2pt] at (#1,0) {\bfseries #2};}
% hito bajo la línea (para alternar y no chocar)
\newcommand{\hitob}[2]{%
  \fill[subjectcolor] (#1,0) circle (3.2pt);
  \node[below, text=ink, align=center, text width=2.6cm, anchor=north,
        yshift=-2pt] at (#1,0) {#2};}

% =====================================================================
%  LÍNEA DE TIEMPO VERTICAL (ideal para listas cronológicas largas)
%  \begin{lineav} \hitov{año}{Título}{detalle} ... \end{lineav}
% =====================================================================
\newcounter{lvrow}
\newenvironment{lineav}{%
  \setcounter{lvrow}{0}%
  \begin{list}{}{\leftmargin=1.6cm \itemsep=0.55em \topsep=0.3em}%
  \renewcommand{\hitov}[3]{\stepcounter{lvrow}%
    \item[\tikz{\node[circle, fill=subjectcolor, text=white, inner sep=0pt,
       minimum size=1.5em, font=\DisplaySB\scriptsize]{\thelvrow};}]%
    {\DisplaySB\color{subjectcolor}##1}\quad{\bfseries ##2}\\[1pt]
    {\small\color{ink!75}##3}}%
}{%
  \end{list}}
\newcommand{\hitov}[3]{}  % placeholder; redefinido dentro del entorno

% =====================================================================
%  MAPA MENTAL — nodo central + ramas radiales
%  \begin{mapa}{NODO CENTRAL}
%     \rama{ángulo}{distancia cm}{texto de la rama}  ...
%  \end{mapa}
% =====================================================================
\newenvironment{mapa}[2][1]{%   % [escala opcional, def=1]{NODO CENTRAL}
  \begin{center}\begin{tikzpicture}[font=\sffamily\small, scale=#1, transform shape]
  \node[circle, fill=subjectcolor, text=white, font=\DisplaySB,
        align=center, text width=2.4cm, inner sep=4pt, minimum size=2.6cm]
        (CENTRO) at (0,0) {#2};
  \renewcommand{\rama}[3]{%
    \node[draw=subjectcolor!55, fill=subjectcolor!8, rounded corners=4pt,
          text=ink, align=center, text width=2.7cm, inner sep=4pt]
          (R##1) at (##1:##2) {##3};
    \draw[subjectcolor!55, line width=1.2pt] (CENTRO) -- (R##1);}%
}{%
  \end{tikzpicture}\end{center}}
\newcommand{\rama}[3]{}  % placeholder; redefinido dentro del entorno
 ; \setsubject{csoc} ; \begin{document}...
% =====================================================================
\documentclass[aspectratio=169, 11pt]{beamer}
\usetheme{default}
\setbeamertemplate{navigation symbols}{}   % sin botones de navegación

\usepackage{fontspec}
\usepackage{polyglossia}\setdefaultlanguage{spanish}
\usepackage{tikz}
\usetikzlibrary{calc, positioning, arrows.meta, decorations.pathreplacing}
\usepackage{amsmath}

% --- TIPOGRAFÍA (idéntica familia que el libro) -----------------------
\setsansfont{Inter}[UprightFont=Inter-Regular, BoldFont=Inter-SemiBold,
  ItalicFont=Inter-Italic, Ligatures=TeX, Numbers=Lining]
\newfontfamily\Display{Inter Display}[Ligatures=TeX]
\newfontfamily\DisplayBlk{Inter Display Black}[Ligatures=TeX]
\newfontfamily\DisplaySB{Inter Display SemiBold}[Ligatures=TeX]
\newfontfamily\DisplayThin{Inter Display Thin}[Ligatures=TeX]

% --- COLOR (mismas 4 paletas que el libro) ----------------------------
\definecolor{mat} {RGB}{ 30,  64, 120}
\definecolor{cnat}{RGB}{ 16, 110,  80}
\definecolor{csoc}{RGB}{165,  75,  45}
\definecolor{eco} {RGB}{ 70,  55, 130}
\definecolor{ink}{RGB}{28,28,30}        % casi-negro para texto
\definecolor{paper}{RGB}{250,250,248}   % blanco cálido de fondo
\colorlet{subjectcolor}{csoc}
% \setsubject re-aplica los colores de beamer además de actualizar el alias
\newcommand{\setsubject}[1]{%
  \colorlet{subjectcolor}{#1}%
  \setbeamercolor{structure}{fg=#1}%
  \setbeamercolor{frametitle}{fg=#1,bg=}%
}

\setbeamercolor{normal text}{fg=ink, bg=paper}
\setbeamercolor{structure}{fg=subjectcolor}
\setbeamerfont{frametitle}{family=\DisplaySB}
\setbeamercolor{frametitle}{fg=subjectcolor, bg=}
\setbeamertemplate{frametitle}{%
  \vspace{0.6em}\hspace{0.2em}%
  {\DisplaySB\large\color{subjectcolor}\insertframetitle}\par
  \vspace{-0.2em}{\color{subjectcolor!30}\rule{\dimexpr\paperwidth-1.4em\relax}{1.2pt}}}

% Pie minimal: instituto · profe · tema · número en cuadro de color
\newcommand{\bmrtema}{}
\newcommand{\bmrprofe}{}   % ← \renewcommand{\bmrprofe}{Prof. Nombre} en cada beamer
\setbeamertemplate{footline}{%
  \hbox{}\vspace{2pt}\hspace{0.6em}%
  {\footnotesize\color{ink!45}ITSTZ%
   \ifx\bmrprofe\empty\else\space·\space\bmrprofe\fi}%
  \hfill{\footnotesize\color{ink!45}\bmrtema}\hspace{0.6em}%
  \tikz[baseline=-3pt]{\node[fill=subjectcolor, text=white, inner xsep=6pt,
     inner ysep=2.5pt, font=\DisplaySB\scriptsize]{\insertframenumber};}%
  \hspace{0.6em}\vspace{4pt}}

% =====================================================================
%  PORTADA y DIVISORES
% =====================================================================
\newcommand{\portada}[3]{% {kicker}{TÍTULO}{subtítulo}
{\setbeamertemplate{footline}{}
\begin{frame}[plain]
\begin{tikzpicture}[remember picture, overlay]
  \fill[paper] (current page.south west) rectangle (current page.north east);
  \node[anchor=north west, text=subjectcolor!12, font=\DisplayBlk]
    at ([xshift=-0.5cm,yshift=1.2cm]current page.north west)
    {\fontsize{190}{190}\selectfont ·};
  \node[anchor=west, text=subjectcolor, font=\DisplaySB\Large]
    at ([xshift=1.2cm, yshift=1.6cm]current page.west) {#1};
  \node[anchor=west, text=ink, font=\DisplayBlk, align=left]
    at ([xshift=1.15cm, yshift=0.2cm]current page.west)
    {\fontsize{40}{44}\selectfont #2};
  \fill[subjectcolor] ([xshift=1.2cm, yshift=-1.4cm]current page.west)
    rectangle ++(3.2cm, 0.16cm);
  \node[anchor=west, text=ink!60, font=\Display, align=left, text width=0.7\paperwidth]
    at ([xshift=1.2cm, yshift=-2.3cm]current page.west) {#3};
\end{tikzpicture}
\end{frame}}}

% Lámina divisoria: solo una frase grande centrada sobre color
\newcommand{\seccion}[1]{%
{\setbeamertemplate{footline}{}
\begin{frame}[plain]
\begin{tikzpicture}[remember picture, overlay]
  \fill[subjectcolor] (current page.south west) rectangle (current page.north east);
  \node[text=white, font=\DisplaySB, align=center, text width=0.85\paperwidth]
    at (current page.center) {\fontsize{30}{36}\selectfont #1};
\end{tikzpicture}
\end{frame}}}

% =====================================================================
%  ELEMENTOS DE CONTENIDO
% =====================================================================
% Una afirmación gigante que llena la lámina
\newcommand{\idea}[1]{%
{\setbeamertemplate{footline}{}
\begin{frame}[plain]
\begin{tikzpicture}[remember picture, overlay]
  \fill[paper] (current page.south west) rectangle (current page.north east);
  \fill[subjectcolor] (current page.north west)
    rectangle ([yshift=-0.32cm]current page.north east);
  \node[text=ink, font=\DisplaySB, align=center,
        text width=0.82\paperwidth]
    at (current page.center)
    {\fontsize{32}{44}\selectfont #1};
\end{tikzpicture}
\end{frame}}}

% Realce inline (fondo color suave, texto oscuro)
\newcommand{\resalta}[1]{\tikz[baseline=(R.base)]{\node[fill=subjectcolor!18,
  rounded corners=2pt, inner xsep=3pt, inner ysep=1pt](R){#1};}}
% Realce fuerte (fondo sólido, texto blanco)
\newcommand{\marca}[1]{\tikz[baseline=(R.base)]{\node[fill=subjectcolor,
  text=white, rounded corners=2pt, inner xsep=4pt, inner ysep=1.5pt](R){\textbf{#1}};}}

% Pregunta viva de apertura
\newcommand{\pregunta}[1]{%
\begin{frame}[plain]
\begin{tikzpicture}[remember picture,overlay]
  \fill[subjectcolor!7] (current page.south west) rectangle (current page.north east);
  \fill[subjectcolor] (current page.south west)
    rectangle ([xshift=0.45cm]current page.north west);
\end{tikzpicture}
\vfill
\hspace{1.5cm}\begin{minipage}{0.78\paperwidth}
  {\color{subjectcolor}\DisplayBlk\fontsize{54}{54}\selectfont ¿}\par
  \vspace{0.25em}
  {\Display\color{ink}\fontsize{22}{28}\selectfont #1}\par
  \vspace{0.4em}
  {\hfill\color{subjectcolor}\DisplayBlk\fontsize{38}{38}\selectfont ?}
\end{minipage}
\vfill
\end{frame}}

% Cierre / acción — label opcional (default: "DE LA QUEJA A LA ACCIÓN")
% Uso: \accion{texto}  o  \accion[ETIQUETA CUSTOM]{texto}
\newcommand{\accion}[2][DE LA QUEJA A LA ACCIÓN]{%
\begin{frame}[plain]
\begin{tikzpicture}[remember picture, overlay]
  \fill[subjectcolor!10] (current page.south west) rectangle (current page.north east);
  \node[anchor=west, text=subjectcolor, font=\DisplaySB]
    at ([xshift=1cm, yshift=2.2cm]current page.west)
    {\large #1};
  \node[anchor=west, text=ink, font=\Display, align=left, text width=0.78\paperwidth]
    at ([xshift=1cm, yshift=0cm]current page.west)
    {\fontsize{20}{26}\selectfont #2};
\end{tikzpicture}
\end{frame}}

% =====================================================================
%  LÍNEA DE TIEMPO HORIZONTAL
%  \begin{lineah}{año_ini}{año_fin}
%     \hito{posición 0-1}{etiqueta}   ...
%  \end{lineah}
% =====================================================================
\newenvironment{lineah}{%
  \begin{center}\begin{tikzpicture}[x=0.92\paperwidth, font=\footnotesize\sffamily]
  \draw[line width=2pt, subjectcolor, -{Stealth[length=5pt]}] (0,0) -- (1.02,0);
}{%
  \end{tikzpicture}\end{center}}
% hito: #1 = posición relativa 0..1 ; #2 = etiqueta (texto corto)
\newcommand{\hito}[2]{%
  \fill[subjectcolor] (#1,0) circle (3.2pt);
  \node[above, text=ink, align=center, text width=2.6cm, anchor=south,
        yshift=2pt] at (#1,0) {\bfseries #2};}
% hito bajo la línea (para alternar y no chocar)
\newcommand{\hitob}[2]{%
  \fill[subjectcolor] (#1,0) circle (3.2pt);
  \node[below, text=ink, align=center, text width=2.6cm, anchor=north,
        yshift=-2pt] at (#1,0) {#2};}

% =====================================================================
%  LÍNEA DE TIEMPO VERTICAL (ideal para listas cronológicas largas)
%  \begin{lineav} \hitov{año}{Título}{detalle} ... \end{lineav}
% =====================================================================
\newcounter{lvrow}
\newenvironment{lineav}{%
  \setcounter{lvrow}{0}%
  \begin{list}{}{\leftmargin=1.6cm \itemsep=0.55em \topsep=0.3em}%
  \renewcommand{\hitov}[3]{\stepcounter{lvrow}%
    \item[\tikz{\node[circle, fill=subjectcolor, text=white, inner sep=0pt,
       minimum size=1.5em, font=\DisplaySB\scriptsize]{\thelvrow};}]%
    {\DisplaySB\color{subjectcolor}##1}\quad{\bfseries ##2}\\[1pt]
    {\small\color{ink!75}##3}}%
}{%
  \end{list}}
\newcommand{\hitov}[3]{}  % placeholder; redefinido dentro del entorno

% =====================================================================
%  MAPA MENTAL — nodo central + ramas radiales
%  \begin{mapa}{NODO CENTRAL}
%     \rama{ángulo}{distancia cm}{texto de la rama}  ...
%  \end{mapa}
% =====================================================================
\newenvironment{mapa}[2][1]{%   % [escala opcional, def=1]{NODO CENTRAL}
  \begin{center}\begin{tikzpicture}[font=\sffamily\small, scale=#1, transform shape]
  \node[circle, fill=subjectcolor, text=white, font=\DisplaySB,
        align=center, text width=2.4cm, inner sep=4pt, minimum size=2.6cm]
        (CENTRO) at (0,0) {#2};
  \renewcommand{\rama}[3]{%
    \node[draw=subjectcolor!55, fill=subjectcolor!8, rounded corners=4pt,
          text=ink, align=center, text width=2.7cm, inner sep=4pt]
          (R##1) at (##1:##2) {##3};
    \draw[subjectcolor!55, line width=1.2pt] (CENTRO) -- (R##1);}%
}{%
  \end{tikzpicture}\end{center}}
\newcommand{\rama}[3]{}  % placeholder; redefinido dentro del entorno

%     % =====================================================================
%  TEMA · MATEMÁTICAS — variante del design system (BEAMER 16:9)
%  Identidad para decks de matemáticas (azul \setsubject{mat}):
%    · FIX del bug de preflight §5/§6: las láminas a sangre del design
%      system (\portada \seccion \idea \pregunta \accion) usan
%      remember picture + current page y salen CORRIDAS hacia arriba con
%      compile_quiet.sh. Aquí se redefinen con `background canvas` +
%      coordenadas absolutas (misma técnica probada del tema eco/cnat).
%    · carga lmodern/colortbl/enumitem (que design/beamer.tex no trae) para
%      math escalable en titulares, \rowcolor en tablas y itemize[clave=valor].
%    · \DisplaySB con BoldFont explícito → \textbf dentro de \idea no avisa.
%    · portada con el personaje "Dr. Cuack" matemático (opcional; se dibuja
%      solo si existe el PNG).
%
%  INVOCACIÓN (tras % =====================================================================
%  DESIGN SYSTEM BEAMER v5 — ITSTZ · 16:9
%  Coherente con el libro: Inter Display (títulos) + Inter (cuerpo),
%  un color de asignatura, fondo claro, texto GRANDE y contrastado.
%  Filosofía slide: una idea por lámina, mucho aire, cero relleno.
%
%  Vocabulario:
%    \portada{kicker}{TÍTULO}{subtítulo}
%    \seccion{texto}                  lámina divisoria a pantalla completa
%    \idea{TEXTO GRANDE}              una afirmación que llena la slide
%    \resalta{texto}                  realce inline (fondo color)
%    frames de contenido normales con \frametitle
%    \begin{lineah} \hito{x}{etiqueta} ... \end{lineah}   timeline horizontal
%    \begin{lineav} \hitov{x}{título}{texto} ... \end{lineav} timeline vertical
%    \begin{mapa}{NODO CENTRAL} \rama{ang}{dist}{texto} ... \end{mapa}
%    \pregunta{texto}                 pregunta viva (apertura)
%    \accion{texto}                   cierre "de la queja a la acción"
%  Uso: \input{design/beamer.tex} ; \setsubject{csoc} ; \begin{document}...
% =====================================================================
\documentclass[aspectratio=169, 11pt]{beamer}
\usetheme{default}
\setbeamertemplate{navigation symbols}{}   % sin botones de navegación

\usepackage{fontspec}
\usepackage{polyglossia}\setdefaultlanguage{spanish}
\usepackage{tikz}
\usetikzlibrary{calc, positioning, arrows.meta, decorations.pathreplacing}
\usepackage{amsmath}

% --- TIPOGRAFÍA (idéntica familia que el libro) -----------------------
\setsansfont{Inter}[UprightFont=Inter-Regular, BoldFont=Inter-SemiBold,
  ItalicFont=Inter-Italic, Ligatures=TeX, Numbers=Lining]
\newfontfamily\Display{Inter Display}[Ligatures=TeX]
\newfontfamily\DisplayBlk{Inter Display Black}[Ligatures=TeX]
\newfontfamily\DisplaySB{Inter Display SemiBold}[Ligatures=TeX]
\newfontfamily\DisplayThin{Inter Display Thin}[Ligatures=TeX]

% --- COLOR (mismas 4 paletas que el libro) ----------------------------
\definecolor{mat} {RGB}{ 30,  64, 120}
\definecolor{cnat}{RGB}{ 16, 110,  80}
\definecolor{csoc}{RGB}{165,  75,  45}
\definecolor{eco} {RGB}{ 70,  55, 130}
\definecolor{ink}{RGB}{28,28,30}        % casi-negro para texto
\definecolor{paper}{RGB}{250,250,248}   % blanco cálido de fondo
\colorlet{subjectcolor}{csoc}
% \setsubject re-aplica los colores de beamer además de actualizar el alias
\newcommand{\setsubject}[1]{%
  \colorlet{subjectcolor}{#1}%
  \setbeamercolor{structure}{fg=#1}%
  \setbeamercolor{frametitle}{fg=#1,bg=}%
}

\setbeamercolor{normal text}{fg=ink, bg=paper}
\setbeamercolor{structure}{fg=subjectcolor}
\setbeamerfont{frametitle}{family=\DisplaySB}
\setbeamercolor{frametitle}{fg=subjectcolor, bg=}
\setbeamertemplate{frametitle}{%
  \vspace{0.6em}\hspace{0.2em}%
  {\DisplaySB\large\color{subjectcolor}\insertframetitle}\par
  \vspace{-0.2em}{\color{subjectcolor!30}\rule{\dimexpr\paperwidth-1.4em\relax}{1.2pt}}}

% Pie minimal: instituto · profe · tema · número en cuadro de color
\newcommand{\bmrtema}{}
\newcommand{\bmrprofe}{}   % ← \renewcommand{\bmrprofe}{Prof. Nombre} en cada beamer
\setbeamertemplate{footline}{%
  \hbox{}\vspace{2pt}\hspace{0.6em}%
  {\footnotesize\color{ink!45}ITSTZ%
   \ifx\bmrprofe\empty\else\space·\space\bmrprofe\fi}%
  \hfill{\footnotesize\color{ink!45}\bmrtema}\hspace{0.6em}%
  \tikz[baseline=-3pt]{\node[fill=subjectcolor, text=white, inner xsep=6pt,
     inner ysep=2.5pt, font=\DisplaySB\scriptsize]{\insertframenumber};}%
  \hspace{0.6em}\vspace{4pt}}

% =====================================================================
%  PORTADA y DIVISORES
% =====================================================================
\newcommand{\portada}[3]{% {kicker}{TÍTULO}{subtítulo}
{\setbeamertemplate{footline}{}
\begin{frame}[plain]
\begin{tikzpicture}[remember picture, overlay]
  \fill[paper] (current page.south west) rectangle (current page.north east);
  \node[anchor=north west, text=subjectcolor!12, font=\DisplayBlk]
    at ([xshift=-0.5cm,yshift=1.2cm]current page.north west)
    {\fontsize{190}{190}\selectfont ·};
  \node[anchor=west, text=subjectcolor, font=\DisplaySB\Large]
    at ([xshift=1.2cm, yshift=1.6cm]current page.west) {#1};
  \node[anchor=west, text=ink, font=\DisplayBlk, align=left]
    at ([xshift=1.15cm, yshift=0.2cm]current page.west)
    {\fontsize{40}{44}\selectfont #2};
  \fill[subjectcolor] ([xshift=1.2cm, yshift=-1.4cm]current page.west)
    rectangle ++(3.2cm, 0.16cm);
  \node[anchor=west, text=ink!60, font=\Display, align=left, text width=0.7\paperwidth]
    at ([xshift=1.2cm, yshift=-2.3cm]current page.west) {#3};
\end{tikzpicture}
\end{frame}}}

% Lámina divisoria: solo una frase grande centrada sobre color
\newcommand{\seccion}[1]{%
{\setbeamertemplate{footline}{}
\begin{frame}[plain]
\begin{tikzpicture}[remember picture, overlay]
  \fill[subjectcolor] (current page.south west) rectangle (current page.north east);
  \node[text=white, font=\DisplaySB, align=center, text width=0.85\paperwidth]
    at (current page.center) {\fontsize{30}{36}\selectfont #1};
\end{tikzpicture}
\end{frame}}}

% =====================================================================
%  ELEMENTOS DE CONTENIDO
% =====================================================================
% Una afirmación gigante que llena la lámina
\newcommand{\idea}[1]{%
{\setbeamertemplate{footline}{}
\begin{frame}[plain]
\begin{tikzpicture}[remember picture, overlay]
  \fill[paper] (current page.south west) rectangle (current page.north east);
  \fill[subjectcolor] (current page.north west)
    rectangle ([yshift=-0.32cm]current page.north east);
  \node[text=ink, font=\DisplaySB, align=center,
        text width=0.82\paperwidth]
    at (current page.center)
    {\fontsize{32}{44}\selectfont #1};
\end{tikzpicture}
\end{frame}}}

% Realce inline (fondo color suave, texto oscuro)
\newcommand{\resalta}[1]{\tikz[baseline=(R.base)]{\node[fill=subjectcolor!18,
  rounded corners=2pt, inner xsep=3pt, inner ysep=1pt](R){#1};}}
% Realce fuerte (fondo sólido, texto blanco)
\newcommand{\marca}[1]{\tikz[baseline=(R.base)]{\node[fill=subjectcolor,
  text=white, rounded corners=2pt, inner xsep=4pt, inner ysep=1.5pt](R){\textbf{#1}};}}

% Pregunta viva de apertura
\newcommand{\pregunta}[1]{%
\begin{frame}[plain]
\begin{tikzpicture}[remember picture,overlay]
  \fill[subjectcolor!7] (current page.south west) rectangle (current page.north east);
  \fill[subjectcolor] (current page.south west)
    rectangle ([xshift=0.45cm]current page.north west);
\end{tikzpicture}
\vfill
\hspace{1.5cm}\begin{minipage}{0.78\paperwidth}
  {\color{subjectcolor}\DisplayBlk\fontsize{54}{54}\selectfont ¿}\par
  \vspace{0.25em}
  {\Display\color{ink}\fontsize{22}{28}\selectfont #1}\par
  \vspace{0.4em}
  {\hfill\color{subjectcolor}\DisplayBlk\fontsize{38}{38}\selectfont ?}
\end{minipage}
\vfill
\end{frame}}

% Cierre / acción — label opcional (default: "DE LA QUEJA A LA ACCIÓN")
% Uso: \accion{texto}  o  \accion[ETIQUETA CUSTOM]{texto}
\newcommand{\accion}[2][DE LA QUEJA A LA ACCIÓN]{%
\begin{frame}[plain]
\begin{tikzpicture}[remember picture, overlay]
  \fill[subjectcolor!10] (current page.south west) rectangle (current page.north east);
  \node[anchor=west, text=subjectcolor, font=\DisplaySB]
    at ([xshift=1cm, yshift=2.2cm]current page.west)
    {\large #1};
  \node[anchor=west, text=ink, font=\Display, align=left, text width=0.78\paperwidth]
    at ([xshift=1cm, yshift=0cm]current page.west)
    {\fontsize{20}{26}\selectfont #2};
\end{tikzpicture}
\end{frame}}

% =====================================================================
%  LÍNEA DE TIEMPO HORIZONTAL
%  \begin{lineah}{año_ini}{año_fin}
%     \hito{posición 0-1}{etiqueta}   ...
%  \end{lineah}
% =====================================================================
\newenvironment{lineah}{%
  \begin{center}\begin{tikzpicture}[x=0.92\paperwidth, font=\footnotesize\sffamily]
  \draw[line width=2pt, subjectcolor, -{Stealth[length=5pt]}] (0,0) -- (1.02,0);
}{%
  \end{tikzpicture}\end{center}}
% hito: #1 = posición relativa 0..1 ; #2 = etiqueta (texto corto)
\newcommand{\hito}[2]{%
  \fill[subjectcolor] (#1,0) circle (3.2pt);
  \node[above, text=ink, align=center, text width=2.6cm, anchor=south,
        yshift=2pt] at (#1,0) {\bfseries #2};}
% hito bajo la línea (para alternar y no chocar)
\newcommand{\hitob}[2]{%
  \fill[subjectcolor] (#1,0) circle (3.2pt);
  \node[below, text=ink, align=center, text width=2.6cm, anchor=north,
        yshift=-2pt] at (#1,0) {#2};}

% =====================================================================
%  LÍNEA DE TIEMPO VERTICAL (ideal para listas cronológicas largas)
%  \begin{lineav} \hitov{año}{Título}{detalle} ... \end{lineav}
% =====================================================================
\newcounter{lvrow}
\newenvironment{lineav}{%
  \setcounter{lvrow}{0}%
  \begin{list}{}{\leftmargin=1.6cm \itemsep=0.55em \topsep=0.3em}%
  \renewcommand{\hitov}[3]{\stepcounter{lvrow}%
    \item[\tikz{\node[circle, fill=subjectcolor, text=white, inner sep=0pt,
       minimum size=1.5em, font=\DisplaySB\scriptsize]{\thelvrow};}]%
    {\DisplaySB\color{subjectcolor}##1}\quad{\bfseries ##2}\\[1pt]
    {\small\color{ink!75}##3}}%
}{%
  \end{list}}
\newcommand{\hitov}[3]{}  % placeholder; redefinido dentro del entorno

% =====================================================================
%  MAPA MENTAL — nodo central + ramas radiales
%  \begin{mapa}{NODO CENTRAL}
%     \rama{ángulo}{distancia cm}{texto de la rama}  ...
%  \end{mapa}
% =====================================================================
\newenvironment{mapa}[2][1]{%   % [escala opcional, def=1]{NODO CENTRAL}
  \begin{center}\begin{tikzpicture}[font=\sffamily\small, scale=#1, transform shape]
  \node[circle, fill=subjectcolor, text=white, font=\DisplaySB,
        align=center, text width=2.4cm, inner sep=4pt, minimum size=2.6cm]
        (CENTRO) at (0,0) {#2};
  \renewcommand{\rama}[3]{%
    \node[draw=subjectcolor!55, fill=subjectcolor!8, rounded corners=4pt,
          text=ink, align=center, text width=2.7cm, inner sep=4pt]
          (R##1) at (##1:##2) {##3};
    \draw[subjectcolor!55, line width=1.2pt] (CENTRO) -- (R##1);}%
}{%
  \end{tikzpicture}\end{center}}
\newcommand{\rama}[3]{}  % placeholder; redefinido dentro del entorno
):
%     % =====================================================================
%  DESIGN SYSTEM BEAMER v5 — ITSTZ · 16:9
%  Coherente con el libro: Inter Display (títulos) + Inter (cuerpo),
%  un color de asignatura, fondo claro, texto GRANDE y contrastado.
%  Filosofía slide: una idea por lámina, mucho aire, cero relleno.
%
%  Vocabulario:
%    \portada{kicker}{TÍTULO}{subtítulo}
%    \seccion{texto}                  lámina divisoria a pantalla completa
%    \idea{TEXTO GRANDE}              una afirmación que llena la slide
%    \resalta{texto}                  realce inline (fondo color)
%    frames de contenido normales con \frametitle
%    \begin{lineah} \hito{x}{etiqueta} ... \end{lineah}   timeline horizontal
%    \begin{lineav} \hitov{x}{título}{texto} ... \end{lineav} timeline vertical
%    \begin{mapa}{NODO CENTRAL} \rama{ang}{dist}{texto} ... \end{mapa}
%    \pregunta{texto}                 pregunta viva (apertura)
%    \accion{texto}                   cierre "de la queja a la acción"
%  Uso: \input{design/beamer.tex} ; \setsubject{csoc} ; \begin{document}...
% =====================================================================
\documentclass[aspectratio=169, 11pt]{beamer}
\usetheme{default}
\setbeamertemplate{navigation symbols}{}   % sin botones de navegación

\usepackage{fontspec}
\usepackage{polyglossia}\setdefaultlanguage{spanish}
\usepackage{tikz}
\usetikzlibrary{calc, positioning, arrows.meta, decorations.pathreplacing}
\usepackage{amsmath}

% --- TIPOGRAFÍA (idéntica familia que el libro) -----------------------
\setsansfont{Inter}[UprightFont=Inter-Regular, BoldFont=Inter-SemiBold,
  ItalicFont=Inter-Italic, Ligatures=TeX, Numbers=Lining]
\newfontfamily\Display{Inter Display}[Ligatures=TeX]
\newfontfamily\DisplayBlk{Inter Display Black}[Ligatures=TeX]
\newfontfamily\DisplaySB{Inter Display SemiBold}[Ligatures=TeX]
\newfontfamily\DisplayThin{Inter Display Thin}[Ligatures=TeX]

% --- COLOR (mismas 4 paletas que el libro) ----------------------------
\definecolor{mat} {RGB}{ 30,  64, 120}
\definecolor{cnat}{RGB}{ 16, 110,  80}
\definecolor{csoc}{RGB}{165,  75,  45}
\definecolor{eco} {RGB}{ 70,  55, 130}
\definecolor{ink}{RGB}{28,28,30}        % casi-negro para texto
\definecolor{paper}{RGB}{250,250,248}   % blanco cálido de fondo
\colorlet{subjectcolor}{csoc}
% \setsubject re-aplica los colores de beamer además de actualizar el alias
\newcommand{\setsubject}[1]{%
  \colorlet{subjectcolor}{#1}%
  \setbeamercolor{structure}{fg=#1}%
  \setbeamercolor{frametitle}{fg=#1,bg=}%
}

\setbeamercolor{normal text}{fg=ink, bg=paper}
\setbeamercolor{structure}{fg=subjectcolor}
\setbeamerfont{frametitle}{family=\DisplaySB}
\setbeamercolor{frametitle}{fg=subjectcolor, bg=}
\setbeamertemplate{frametitle}{%
  \vspace{0.6em}\hspace{0.2em}%
  {\DisplaySB\large\color{subjectcolor}\insertframetitle}\par
  \vspace{-0.2em}{\color{subjectcolor!30}\rule{\dimexpr\paperwidth-1.4em\relax}{1.2pt}}}

% Pie minimal: instituto · profe · tema · número en cuadro de color
\newcommand{\bmrtema}{}
\newcommand{\bmrprofe}{}   % ← \renewcommand{\bmrprofe}{Prof. Nombre} en cada beamer
\setbeamertemplate{footline}{%
  \hbox{}\vspace{2pt}\hspace{0.6em}%
  {\footnotesize\color{ink!45}ITSTZ%
   \ifx\bmrprofe\empty\else\space·\space\bmrprofe\fi}%
  \hfill{\footnotesize\color{ink!45}\bmrtema}\hspace{0.6em}%
  \tikz[baseline=-3pt]{\node[fill=subjectcolor, text=white, inner xsep=6pt,
     inner ysep=2.5pt, font=\DisplaySB\scriptsize]{\insertframenumber};}%
  \hspace{0.6em}\vspace{4pt}}

% =====================================================================
%  PORTADA y DIVISORES
% =====================================================================
\newcommand{\portada}[3]{% {kicker}{TÍTULO}{subtítulo}
{\setbeamertemplate{footline}{}
\begin{frame}[plain]
\begin{tikzpicture}[remember picture, overlay]
  \fill[paper] (current page.south west) rectangle (current page.north east);
  \node[anchor=north west, text=subjectcolor!12, font=\DisplayBlk]
    at ([xshift=-0.5cm,yshift=1.2cm]current page.north west)
    {\fontsize{190}{190}\selectfont ·};
  \node[anchor=west, text=subjectcolor, font=\DisplaySB\Large]
    at ([xshift=1.2cm, yshift=1.6cm]current page.west) {#1};
  \node[anchor=west, text=ink, font=\DisplayBlk, align=left]
    at ([xshift=1.15cm, yshift=0.2cm]current page.west)
    {\fontsize{40}{44}\selectfont #2};
  \fill[subjectcolor] ([xshift=1.2cm, yshift=-1.4cm]current page.west)
    rectangle ++(3.2cm, 0.16cm);
  \node[anchor=west, text=ink!60, font=\Display, align=left, text width=0.7\paperwidth]
    at ([xshift=1.2cm, yshift=-2.3cm]current page.west) {#3};
\end{tikzpicture}
\end{frame}}}

% Lámina divisoria: solo una frase grande centrada sobre color
\newcommand{\seccion}[1]{%
{\setbeamertemplate{footline}{}
\begin{frame}[plain]
\begin{tikzpicture}[remember picture, overlay]
  \fill[subjectcolor] (current page.south west) rectangle (current page.north east);
  \node[text=white, font=\DisplaySB, align=center, text width=0.85\paperwidth]
    at (current page.center) {\fontsize{30}{36}\selectfont #1};
\end{tikzpicture}
\end{frame}}}

% =====================================================================
%  ELEMENTOS DE CONTENIDO
% =====================================================================
% Una afirmación gigante que llena la lámina
\newcommand{\idea}[1]{%
{\setbeamertemplate{footline}{}
\begin{frame}[plain]
\begin{tikzpicture}[remember picture, overlay]
  \fill[paper] (current page.south west) rectangle (current page.north east);
  \fill[subjectcolor] (current page.north west)
    rectangle ([yshift=-0.32cm]current page.north east);
  \node[text=ink, font=\DisplaySB, align=center,
        text width=0.82\paperwidth]
    at (current page.center)
    {\fontsize{32}{44}\selectfont #1};
\end{tikzpicture}
\end{frame}}}

% Realce inline (fondo color suave, texto oscuro)
\newcommand{\resalta}[1]{\tikz[baseline=(R.base)]{\node[fill=subjectcolor!18,
  rounded corners=2pt, inner xsep=3pt, inner ysep=1pt](R){#1};}}
% Realce fuerte (fondo sólido, texto blanco)
\newcommand{\marca}[1]{\tikz[baseline=(R.base)]{\node[fill=subjectcolor,
  text=white, rounded corners=2pt, inner xsep=4pt, inner ysep=1.5pt](R){\textbf{#1}};}}

% Pregunta viva de apertura
\newcommand{\pregunta}[1]{%
\begin{frame}[plain]
\begin{tikzpicture}[remember picture,overlay]
  \fill[subjectcolor!7] (current page.south west) rectangle (current page.north east);
  \fill[subjectcolor] (current page.south west)
    rectangle ([xshift=0.45cm]current page.north west);
\end{tikzpicture}
\vfill
\hspace{1.5cm}\begin{minipage}{0.78\paperwidth}
  {\color{subjectcolor}\DisplayBlk\fontsize{54}{54}\selectfont ¿}\par
  \vspace{0.25em}
  {\Display\color{ink}\fontsize{22}{28}\selectfont #1}\par
  \vspace{0.4em}
  {\hfill\color{subjectcolor}\DisplayBlk\fontsize{38}{38}\selectfont ?}
\end{minipage}
\vfill
\end{frame}}

% Cierre / acción — label opcional (default: "DE LA QUEJA A LA ACCIÓN")
% Uso: \accion{texto}  o  \accion[ETIQUETA CUSTOM]{texto}
\newcommand{\accion}[2][DE LA QUEJA A LA ACCIÓN]{%
\begin{frame}[plain]
\begin{tikzpicture}[remember picture, overlay]
  \fill[subjectcolor!10] (current page.south west) rectangle (current page.north east);
  \node[anchor=west, text=subjectcolor, font=\DisplaySB]
    at ([xshift=1cm, yshift=2.2cm]current page.west)
    {\large #1};
  \node[anchor=west, text=ink, font=\Display, align=left, text width=0.78\paperwidth]
    at ([xshift=1cm, yshift=0cm]current page.west)
    {\fontsize{20}{26}\selectfont #2};
\end{tikzpicture}
\end{frame}}

% =====================================================================
%  LÍNEA DE TIEMPO HORIZONTAL
%  \begin{lineah}{año_ini}{año_fin}
%     \hito{posición 0-1}{etiqueta}   ...
%  \end{lineah}
% =====================================================================
\newenvironment{lineah}{%
  \begin{center}\begin{tikzpicture}[x=0.92\paperwidth, font=\footnotesize\sffamily]
  \draw[line width=2pt, subjectcolor, -{Stealth[length=5pt]}] (0,0) -- (1.02,0);
}{%
  \end{tikzpicture}\end{center}}
% hito: #1 = posición relativa 0..1 ; #2 = etiqueta (texto corto)
\newcommand{\hito}[2]{%
  \fill[subjectcolor] (#1,0) circle (3.2pt);
  \node[above, text=ink, align=center, text width=2.6cm, anchor=south,
        yshift=2pt] at (#1,0) {\bfseries #2};}
% hito bajo la línea (para alternar y no chocar)
\newcommand{\hitob}[2]{%
  \fill[subjectcolor] (#1,0) circle (3.2pt);
  \node[below, text=ink, align=center, text width=2.6cm, anchor=north,
        yshift=-2pt] at (#1,0) {#2};}

% =====================================================================
%  LÍNEA DE TIEMPO VERTICAL (ideal para listas cronológicas largas)
%  \begin{lineav} \hitov{año}{Título}{detalle} ... \end{lineav}
% =====================================================================
\newcounter{lvrow}
\newenvironment{lineav}{%
  \setcounter{lvrow}{0}%
  \begin{list}{}{\leftmargin=1.6cm \itemsep=0.55em \topsep=0.3em}%
  \renewcommand{\hitov}[3]{\stepcounter{lvrow}%
    \item[\tikz{\node[circle, fill=subjectcolor, text=white, inner sep=0pt,
       minimum size=1.5em, font=\DisplaySB\scriptsize]{\thelvrow};}]%
    {\DisplaySB\color{subjectcolor}##1}\quad{\bfseries ##2}\\[1pt]
    {\small\color{ink!75}##3}}%
}{%
  \end{list}}
\newcommand{\hitov}[3]{}  % placeholder; redefinido dentro del entorno

% =====================================================================
%  MAPA MENTAL — nodo central + ramas radiales
%  \begin{mapa}{NODO CENTRAL}
%     \rama{ángulo}{distancia cm}{texto de la rama}  ...
%  \end{mapa}
% =====================================================================
\newenvironment{mapa}[2][1]{%   % [escala opcional, def=1]{NODO CENTRAL}
  \begin{center}\begin{tikzpicture}[font=\sffamily\small, scale=#1, transform shape]
  \node[circle, fill=subjectcolor, text=white, font=\DisplaySB,
        align=center, text width=2.4cm, inner sep=4pt, minimum size=2.6cm]
        (CENTRO) at (0,0) {#2};
  \renewcommand{\rama}[3]{%
    \node[draw=subjectcolor!55, fill=subjectcolor!8, rounded corners=4pt,
          text=ink, align=center, text width=2.7cm, inner sep=4pt]
          (R##1) at (##1:##2) {##3};
    \draw[subjectcolor!55, line width=1.2pt] (CENTRO) -- (R##1);}%
}{%
  \end{tikzpicture}\end{center}}
\newcommand{\rama}[3]{}  % placeholder; redefinido dentro del entorno

%     % =====================================================================
%  TEMA · MATEMÁTICAS — variante del design system (BEAMER 16:9)
%  Identidad para decks de matemáticas (azul \setsubject{mat}):
%    · FIX del bug de preflight §5/§6: las láminas a sangre del design
%      system (\portada \seccion \idea \pregunta \accion) usan
%      remember picture + current page y salen CORRIDAS hacia arriba con
%      compile_quiet.sh. Aquí se redefinen con `background canvas` +
%      coordenadas absolutas (misma técnica probada del tema eco/cnat).
%    · carga lmodern/colortbl/enumitem (que design/beamer.tex no trae) para
%      math escalable en titulares, \rowcolor en tablas y itemize[clave=valor].
%    · \DisplaySB con BoldFont explícito → \textbf dentro de \idea no avisa.
%    · portada con el personaje "Dr. Cuack" matemático (opcional; se dibuja
%      solo si existe el PNG).
%
%  INVOCACIÓN (tras \input{design/beamer.tex}):
%     \input{design/beamer.tex}
%     \input{design/tema-mat-beamer.tex}
% =====================================================================

\usepackage{fontawesome5}         % \faLightbulb \faCheckCircle ...
\usepackage{lmodern}              % math Latin Modern ESCALABLE (titulares con math)
\usepackage{colortbl}             % \rowcolor con tintes !NN (beamer no trae colortbl)
\usepackage{enumitem}             % itemize/enumerate[clave=valor]
\setlist{leftmargin=1.4em, itemsep=3pt, topsep=2pt}

% \DisplaySB (peso, no familia) no autodetecta BOLD → \textbf avisa. BoldFont explícito.
\renewfontfamily\DisplaySB{Inter Display SemiBold}[Ligatures=TeX,
  BoldFont={Inter Display SemiBold}]

\setsubject{mat}                  % azul RGB(30,64,120)

% Titulares grandes NO deben partir palabras. Dentro de nodos TikZ con `text
% width`, subir \hyphenpenalty NO basta: se anula el guion del font ACTIVO con
% \hyphenchar\font=-1 DESPUÉS del \selectfont del tamaño usado → \nh.
\newcommand{\nohyf}{\hyphenpenalty=10000\relax\exhyphenpenalty=10000\relax}
\newcommand{\nh}{\hyphenchar\font=-1\relax}

% Ruta al personaje (se puede sobreescribir por deck con \renewcommand)
\newcommand{\drcuack}{assets/personaje/calcado/png/dr-cuack-matematico.png}

% =====================================================================
%  LÁMINAS A SANGRE — redefinidas con background canvas (fix §5/§6)
% =====================================================================

% --- PORTADA: texto a la izquierda + Dr. Cuack a la derecha -----------
\renewcommand{\portada}[3]{% {kicker}{TÍTULO}{subtítulo}
{\setbeamertemplate{footline}{}
\setbeamertemplate{background canvas}{%
  \begin{tikzpicture}
    \useasboundingbox (0,0) rectangle (\paperwidth,\paperheight);
    \fill[paper] (0,0) rectangle (\paperwidth,\paperheight);
    \IfFileExists{\drcuack}{%
      \node[anchor=east] at (\paperwidth-0.7,\paperheight*0.46)
        {\includegraphics[height=0.62\paperheight]{\drcuack}};}{}%
    \node[anchor=west, text=ink, align=left, text width=0.56\paperwidth]
      at (1.15,\paperheight*0.5) {\nohyf%
        {\DisplaySB\Large\nh\color{subjectcolor} #1}\\[0.5em]
        {\DisplayBlk\fontsize{36}{40}\selectfont\nh #2}\\[0.7em]
        {\color{subjectcolor}\rule{3.2cm}{0.16cm}}\\[0.6em]
        {\Display\nh\color{ink!60} #3}};
  \end{tikzpicture}}
\begin{frame}[plain]\strut\end{frame}}}

% --- SECCIÓN (divisor): frase blanca centrada sobre azul ---------------
\renewcommand{\seccion}[1]{%
{\setbeamertemplate{footline}{}
\setbeamertemplate{background canvas}{%
  \begin{tikzpicture}
    \useasboundingbox (0,0) rectangle (\paperwidth,\paperheight);
    \fill[subjectcolor] (0,0) rectangle (\paperwidth,\paperheight);
    \node[text=white!20, anchor=south west]
      at (0.4,0.3) {\fontsize{110}{110}\selectfont\(\times\)};
    \node[text=white, font=\DisplaySB, align=center, text width=0.85\paperwidth]
      at (\paperwidth*0.5,\paperheight*0.5) {\nohyf\fontsize{28}{36}\selectfont\nh #1};
  \end{tikzpicture}}
\begin{frame}[plain]\strut\end{frame}}}

% --- IDEA: afirmación gigante con filete azul superior ----------------
\renewcommand{\idea}[1]{%
{\setbeamertemplate{footline}{}
\setbeamertemplate{background canvas}{%
  \begin{tikzpicture}
    \useasboundingbox (0,0) rectangle (\paperwidth,\paperheight);
    \fill[paper] (0,0) rectangle (\paperwidth,\paperheight);
    \fill[subjectcolor] (0,\paperheight-0.32) rectangle (\paperwidth,\paperheight);
    \node[text=ink, font=\DisplaySB, align=center, text width=0.82\paperwidth]
      at (\paperwidth*0.5,\paperheight*0.5) {\nohyf\fontsize{28}{44}\selectfont\nh #1};
  \end{tikzpicture}}
\begin{frame}[plain]\strut\end{frame}}}

% --- PREGUNTA VIVA: barra azul lateral + signos de interrogación -------
\renewcommand{\pregunta}[1]{%
{\setbeamertemplate{footline}{}
\setbeamertemplate{background canvas}{%
  \begin{tikzpicture}
    \useasboundingbox (0,0) rectangle (\paperwidth,\paperheight);
    \fill[subjectcolor!7] (0,0) rectangle (\paperwidth,\paperheight);
    \fill[subjectcolor] (0,0) rectangle (0.42,\paperheight);
    \node[anchor=west, text=ink, align=left, text width=0.74\paperwidth]
      at (1.5,\paperheight*0.5) {\nohyf%
        {\color{subjectcolor}\DisplayBlk\fontsize{40}{40}\selectfont ¿}\\[0.1em]
        {\Display\fontsize{22}{34}\selectfont\nh #1}\\[0.15em]
        {\color{subjectcolor}\DisplayBlk\fontsize{30}{30}\selectfont ?}};
  \end{tikzpicture}}
\begin{frame}[plain]\strut\end{frame}}}

% --- ACCIÓN / cierre --------------------------------------------------
\renewcommand{\accion}[2][EN CORTO]{%
{\setbeamertemplate{background canvas}{%
  \begin{tikzpicture}
    \useasboundingbox (0,0) rectangle (\paperwidth,\paperheight);
    \fill[subjectcolor!10] (0,0) rectangle (\paperwidth,\paperheight);
    \node[anchor=west, align=left, text width=0.82\paperwidth]
      at (1,\paperheight*0.5) {\nohyf%
        {\DisplaySB\large\nh\color{subjectcolor}\faCheckCircle\enspace #1}\\[0.7em]
        {\nohyf\Display\fontsize{19}{26}\selectfont\nh\color{ink} #2}};
  \end{tikzpicture}}
\begin{frame}[plain]\strut\end{frame}}}

% =====================================================================

\usepackage{fontawesome5}         % \faLightbulb \faCheckCircle ...
\usepackage{lmodern}              % math Latin Modern ESCALABLE (titulares con math)
\usepackage{colortbl}             % \rowcolor con tintes !NN (beamer no trae colortbl)
\usepackage{enumitem}             % itemize/enumerate[clave=valor]
\setlist{leftmargin=1.4em, itemsep=3pt, topsep=2pt}

% \DisplaySB (peso, no familia) no autodetecta BOLD → \textbf avisa. BoldFont explícito.
\renewfontfamily\DisplaySB{Inter Display SemiBold}[Ligatures=TeX,
  BoldFont={Inter Display SemiBold}]

\setsubject{mat}                  % azul RGB(30,64,120)

% Titulares grandes NO deben partir palabras. Dentro de nodos TikZ con `text
% width`, subir \hyphenpenalty NO basta: se anula el guion del font ACTIVO con
% \hyphenchar\font=-1 DESPUÉS del \selectfont del tamaño usado → \nh.
\newcommand{\nohyf}{\hyphenpenalty=10000\relax\exhyphenpenalty=10000\relax}
\newcommand{\nh}{\hyphenchar\font=-1\relax}

% Ruta al personaje (se puede sobreescribir por deck con \renewcommand)
\newcommand{\drcuack}{assets/personaje/calcado/png/dr-cuack-matematico.png}

% =====================================================================
%  LÁMINAS A SANGRE — redefinidas con background canvas (fix §5/§6)
% =====================================================================

% --- PORTADA: texto a la izquierda + Dr. Cuack a la derecha -----------
\renewcommand{\portada}[3]{% {kicker}{TÍTULO}{subtítulo}
{\setbeamertemplate{footline}{}
\setbeamertemplate{background canvas}{%
  \begin{tikzpicture}
    \useasboundingbox (0,0) rectangle (\paperwidth,\paperheight);
    \fill[paper] (0,0) rectangle (\paperwidth,\paperheight);
    \IfFileExists{\drcuack}{%
      \node[anchor=east] at (\paperwidth-0.7,\paperheight*0.46)
        {\includegraphics[height=0.62\paperheight]{\drcuack}};}{}%
    \node[anchor=west, text=ink, align=left, text width=0.56\paperwidth]
      at (1.15,\paperheight*0.5) {\nohyf%
        {\DisplaySB\Large\nh\color{subjectcolor} #1}\\[0.5em]
        {\DisplayBlk\fontsize{36}{40}\selectfont\nh #2}\\[0.7em]
        {\color{subjectcolor}\rule{3.2cm}{0.16cm}}\\[0.6em]
        {\Display\nh\color{ink!60} #3}};
  \end{tikzpicture}}
\begin{frame}[plain]\strut\end{frame}}}

% --- SECCIÓN (divisor): frase blanca centrada sobre azul ---------------
\renewcommand{\seccion}[1]{%
{\setbeamertemplate{footline}{}
\setbeamertemplate{background canvas}{%
  \begin{tikzpicture}
    \useasboundingbox (0,0) rectangle (\paperwidth,\paperheight);
    \fill[subjectcolor] (0,0) rectangle (\paperwidth,\paperheight);
    \node[text=white!20, anchor=south west]
      at (0.4,0.3) {\fontsize{110}{110}\selectfont\(\times\)};
    \node[text=white, font=\DisplaySB, align=center, text width=0.85\paperwidth]
      at (\paperwidth*0.5,\paperheight*0.5) {\nohyf\fontsize{28}{36}\selectfont\nh #1};
  \end{tikzpicture}}
\begin{frame}[plain]\strut\end{frame}}}

% --- IDEA: afirmación gigante con filete azul superior ----------------
\renewcommand{\idea}[1]{%
{\setbeamertemplate{footline}{}
\setbeamertemplate{background canvas}{%
  \begin{tikzpicture}
    \useasboundingbox (0,0) rectangle (\paperwidth,\paperheight);
    \fill[paper] (0,0) rectangle (\paperwidth,\paperheight);
    \fill[subjectcolor] (0,\paperheight-0.32) rectangle (\paperwidth,\paperheight);
    \node[text=ink, font=\DisplaySB, align=center, text width=0.82\paperwidth]
      at (\paperwidth*0.5,\paperheight*0.5) {\nohyf\fontsize{28}{44}\selectfont\nh #1};
  \end{tikzpicture}}
\begin{frame}[plain]\strut\end{frame}}}

% --- PREGUNTA VIVA: barra azul lateral + signos de interrogación -------
\renewcommand{\pregunta}[1]{%
{\setbeamertemplate{footline}{}
\setbeamertemplate{background canvas}{%
  \begin{tikzpicture}
    \useasboundingbox (0,0) rectangle (\paperwidth,\paperheight);
    \fill[subjectcolor!7] (0,0) rectangle (\paperwidth,\paperheight);
    \fill[subjectcolor] (0,0) rectangle (0.42,\paperheight);
    \node[anchor=west, text=ink, align=left, text width=0.74\paperwidth]
      at (1.5,\paperheight*0.5) {\nohyf%
        {\color{subjectcolor}\DisplayBlk\fontsize{40}{40}\selectfont ¿}\\[0.1em]
        {\Display\fontsize{22}{34}\selectfont\nh #1}\\[0.15em]
        {\color{subjectcolor}\DisplayBlk\fontsize{30}{30}\selectfont ?}};
  \end{tikzpicture}}
\begin{frame}[plain]\strut\end{frame}}}

% --- ACCIÓN / cierre --------------------------------------------------
\renewcommand{\accion}[2][EN CORTO]{%
{\setbeamertemplate{background canvas}{%
  \begin{tikzpicture}
    \useasboundingbox (0,0) rectangle (\paperwidth,\paperheight);
    \fill[subjectcolor!10] (0,0) rectangle (\paperwidth,\paperheight);
    \node[anchor=west, align=left, text width=0.82\paperwidth]
      at (1,\paperheight*0.5) {\nohyf%
        {\DisplaySB\large\nh\color{subjectcolor}\faCheckCircle\enspace #1}\\[0.7em]
        {\nohyf\Display\fontsize{19}{26}\selectfont\nh\color{ink} #2}};
  \end{tikzpicture}}
\begin{frame}[plain]\strut\end{frame}}}

% =====================================================================

\usepackage{fontawesome5}         % \faLightbulb \faCheckCircle ...
\usepackage{lmodern}              % math Latin Modern ESCALABLE (titulares con math)
\usepackage{colortbl}             % \rowcolor con tintes !NN (beamer no trae colortbl)
\usepackage{enumitem}             % itemize/enumerate[clave=valor]
\setlist{leftmargin=1.4em, itemsep=3pt, topsep=2pt}

% \DisplaySB (peso, no familia) no autodetecta BOLD → \textbf avisa. BoldFont explícito.
\renewfontfamily\DisplaySB{Inter Display SemiBold}[Ligatures=TeX,
  BoldFont={Inter Display SemiBold}]

\setsubject{mat}                  % azul RGB(30,64,120)

% Titulares grandes NO deben partir palabras. Dentro de nodos TikZ con `text
% width`, subir \hyphenpenalty NO basta: se anula el guion del font ACTIVO con
% \hyphenchar\font=-1 DESPUÉS del \selectfont del tamaño usado → \nh.
\newcommand{\nohyf}{\hyphenpenalty=10000\relax\exhyphenpenalty=10000\relax}
\newcommand{\nh}{\hyphenchar\font=-1\relax}

% Ruta al personaje (se puede sobreescribir por deck con \renewcommand)
\newcommand{\drcuack}{assets/personaje/calcado/png/dr-cuack-matematico.png}

% =====================================================================
%  LÁMINAS A SANGRE — redefinidas con background canvas (fix §5/§6)
% =====================================================================

% --- PORTADA: texto a la izquierda + Dr. Cuack a la derecha -----------
\renewcommand{\portada}[3]{% {kicker}{TÍTULO}{subtítulo}
{\setbeamertemplate{footline}{}
\setbeamertemplate{background canvas}{%
  \begin{tikzpicture}
    \useasboundingbox (0,0) rectangle (\paperwidth,\paperheight);
    \fill[paper] (0,0) rectangle (\paperwidth,\paperheight);
    \IfFileExists{\drcuack}{%
      \node[anchor=east] at (\paperwidth-0.7,\paperheight*0.46)
        {\includegraphics[height=0.62\paperheight]{\drcuack}};}{}%
    \node[anchor=west, text=ink, align=left, text width=0.56\paperwidth]
      at (1.15,\paperheight*0.5) {\nohyf%
        {\DisplaySB\Large\nh\color{subjectcolor} #1}\\[0.5em]
        {\DisplayBlk\fontsize{36}{40}\selectfont\nh #2}\\[0.7em]
        {\color{subjectcolor}\rule{3.2cm}{0.16cm}}\\[0.6em]
        {\Display\nh\color{ink!60} #3}};
  \end{tikzpicture}}
\begin{frame}[plain]\strut\end{frame}}}

% --- SECCIÓN (divisor): frase blanca centrada sobre azul ---------------
\renewcommand{\seccion}[1]{%
{\setbeamertemplate{footline}{}
\setbeamertemplate{background canvas}{%
  \begin{tikzpicture}
    \useasboundingbox (0,0) rectangle (\paperwidth,\paperheight);
    \fill[subjectcolor] (0,0) rectangle (\paperwidth,\paperheight);
    \node[text=white!20, anchor=south west]
      at (0.4,0.3) {\fontsize{110}{110}\selectfont\(\times\)};
    \node[text=white, font=\DisplaySB, align=center, text width=0.85\paperwidth]
      at (\paperwidth*0.5,\paperheight*0.5) {\nohyf\fontsize{28}{36}\selectfont\nh #1};
  \end{tikzpicture}}
\begin{frame}[plain]\strut\end{frame}}}

% --- IDEA: afirmación gigante con filete azul superior ----------------
\renewcommand{\idea}[1]{%
{\setbeamertemplate{footline}{}
\setbeamertemplate{background canvas}{%
  \begin{tikzpicture}
    \useasboundingbox (0,0) rectangle (\paperwidth,\paperheight);
    \fill[paper] (0,0) rectangle (\paperwidth,\paperheight);
    \fill[subjectcolor] (0,\paperheight-0.32) rectangle (\paperwidth,\paperheight);
    \node[text=ink, font=\DisplaySB, align=center, text width=0.82\paperwidth]
      at (\paperwidth*0.5,\paperheight*0.5) {\nohyf\fontsize{28}{44}\selectfont\nh #1};
  \end{tikzpicture}}
\begin{frame}[plain]\strut\end{frame}}}

% --- PREGUNTA VIVA: barra azul lateral + signos de interrogación -------
\renewcommand{\pregunta}[1]{%
{\setbeamertemplate{footline}{}
\setbeamertemplate{background canvas}{%
  \begin{tikzpicture}
    \useasboundingbox (0,0) rectangle (\paperwidth,\paperheight);
    \fill[subjectcolor!7] (0,0) rectangle (\paperwidth,\paperheight);
    \fill[subjectcolor] (0,0) rectangle (0.42,\paperheight);
    \node[anchor=west, text=ink, align=left, text width=0.74\paperwidth]
      at (1.5,\paperheight*0.5) {\nohyf%
        {\color{subjectcolor}\DisplayBlk\fontsize{40}{40}\selectfont ¿}\\[0.1em]
        {\Display\fontsize{22}{34}\selectfont\nh #1}\\[0.15em]
        {\color{subjectcolor}\DisplayBlk\fontsize{30}{30}\selectfont ?}};
  \end{tikzpicture}}
\begin{frame}[plain]\strut\end{frame}}}

% --- ACCIÓN / cierre --------------------------------------------------
\renewcommand{\accion}[2][EN CORTO]{%
{\setbeamertemplate{background canvas}{%
  \begin{tikzpicture}
    \useasboundingbox (0,0) rectangle (\paperwidth,\paperheight);
    \fill[subjectcolor!10] (0,0) rectangle (\paperwidth,\paperheight);
    \node[anchor=west, align=left, text width=0.82\paperwidth]
      at (1,\paperheight*0.5) {\nohyf%
        {\DisplaySB\large\nh\color{subjectcolor}\faCheckCircle\enspace #1}\\[0.7em]
        {\nohyf\Display\fontsize{19}{26}\selectfont\nh\color{ink} #2}};
  \end{tikzpicture}}
\begin{frame}[plain]\strut\end{frame}}}

% =====================================================================

\usepackage{fontawesome5}         % \faLightbulb \faCheckCircle ...
\usepackage{lmodern}              % math Latin Modern ESCALABLE (titulares con math)
\usepackage{colortbl}             % \rowcolor con tintes !NN (beamer no trae colortbl)
\usepackage{enumitem}             % itemize/enumerate[clave=valor]
\setlist{leftmargin=1.4em, itemsep=3pt, topsep=2pt}

% \DisplaySB (peso, no familia) no autodetecta BOLD → \textbf avisa. BoldFont explícito.
\renewfontfamily\DisplaySB{Inter Display SemiBold}[Ligatures=TeX,
  BoldFont={Inter Display SemiBold}]

\setsubject{mat}                  % azul RGB(30,64,120)

% Titulares grandes NO deben partir palabras. Dentro de nodos TikZ con `text
% width`, subir \hyphenpenalty NO basta: se anula el guion del font ACTIVO con
% \hyphenchar\font=-1 DESPUÉS del \selectfont del tamaño usado → \nh.
\newcommand{\nohyf}{\hyphenpenalty=10000\relax\exhyphenpenalty=10000\relax}
\newcommand{\nh}{\hyphenchar\font=-1\relax}

% Ruta al personaje (se puede sobreescribir por deck con \renewcommand)
\newcommand{\drcuack}{assets/personaje/calcado/png/dr-cuack-matematico.png}

% =====================================================================
%  LÁMINAS A SANGRE — redefinidas con background canvas (fix §5/§6)
% =====================================================================

% --- PORTADA: texto a la izquierda + Dr. Cuack a la derecha -----------
\renewcommand{\portada}[3]{% {kicker}{TÍTULO}{subtítulo}
{\setbeamertemplate{footline}{}
\setbeamertemplate{background canvas}{%
  \begin{tikzpicture}
    \useasboundingbox (0,0) rectangle (\paperwidth,\paperheight);
    \fill[paper] (0,0) rectangle (\paperwidth,\paperheight);
    \IfFileExists{\drcuack}{%
      \node[anchor=east] at (\paperwidth-0.7,\paperheight*0.46)
        {\includegraphics[height=0.62\paperheight]{\drcuack}};}{}%
    \node[anchor=west, text=ink, align=left, text width=0.56\paperwidth]
      at (1.15,\paperheight*0.5) {\nohyf%
        {\DisplaySB\Large\nh\color{subjectcolor} #1}\\[0.5em]
        {\DisplayBlk\fontsize{36}{40}\selectfont\nh #2}\\[0.7em]
        {\color{subjectcolor}\rule{3.2cm}{0.16cm}}\\[0.6em]
        {\Display\nh\color{ink!60} #3}};
  \end{tikzpicture}}
\begin{frame}[plain]\strut\end{frame}}}

% --- SECCIÓN (divisor): frase blanca centrada sobre azul ---------------
\renewcommand{\seccion}[1]{%
{\setbeamertemplate{footline}{}
\setbeamertemplate{background canvas}{%
  \begin{tikzpicture}
    \useasboundingbox (0,0) rectangle (\paperwidth,\paperheight);
    \fill[subjectcolor] (0,0) rectangle (\paperwidth,\paperheight);
    \node[text=white!20, anchor=south west]
      at (0.4,0.3) {\fontsize{110}{110}\selectfont\(\times\)};
    \node[text=white, font=\DisplaySB, align=center, text width=0.85\paperwidth]
      at (\paperwidth*0.5,\paperheight*0.5) {\nohyf\fontsize{28}{36}\selectfont\nh #1};
  \end{tikzpicture}}
\begin{frame}[plain]\strut\end{frame}}}

% --- IDEA: afirmación gigante con filete azul superior ----------------
\renewcommand{\idea}[1]{%
{\setbeamertemplate{footline}{}
\setbeamertemplate{background canvas}{%
  \begin{tikzpicture}
    \useasboundingbox (0,0) rectangle (\paperwidth,\paperheight);
    \fill[paper] (0,0) rectangle (\paperwidth,\paperheight);
    \fill[subjectcolor] (0,\paperheight-0.32) rectangle (\paperwidth,\paperheight);
    \node[text=ink, font=\DisplaySB, align=center, text width=0.82\paperwidth]
      at (\paperwidth*0.5,\paperheight*0.5) {\nohyf\fontsize{28}{44}\selectfont\nh #1};
  \end{tikzpicture}}
\begin{frame}[plain]\strut\end{frame}}}

% --- PREGUNTA VIVA: barra azul lateral + signos de interrogación -------
\renewcommand{\pregunta}[1]{%
{\setbeamertemplate{footline}{}
\setbeamertemplate{background canvas}{%
  \begin{tikzpicture}
    \useasboundingbox (0,0) rectangle (\paperwidth,\paperheight);
    \fill[subjectcolor!7] (0,0) rectangle (\paperwidth,\paperheight);
    \fill[subjectcolor] (0,0) rectangle (0.42,\paperheight);
    \node[anchor=west, text=ink, align=left, text width=0.74\paperwidth]
      at (1.5,\paperheight*0.5) {\nohyf%
        {\color{subjectcolor}\DisplayBlk\fontsize{40}{40}\selectfont ¿}\\[0.1em]
        {\Display\fontsize{22}{34}\selectfont\nh #1}\\[0.15em]
        {\color{subjectcolor}\DisplayBlk\fontsize{30}{30}\selectfont ?}};
  \end{tikzpicture}}
\begin{frame}[plain]\strut\end{frame}}}

% --- ACCIÓN / cierre --------------------------------------------------
\renewcommand{\accion}[2][EN CORTO]{%
{\setbeamertemplate{background canvas}{%
  \begin{tikzpicture}
    \useasboundingbox (0,0) rectangle (\paperwidth,\paperheight);
    \fill[subjectcolor!10] (0,0) rectangle (\paperwidth,\paperheight);
    \node[anchor=west, align=left, text width=0.82\paperwidth]
      at (1,\paperheight*0.5) {\nohyf%
        {\DisplaySB\large\nh\color{subjectcolor}\faCheckCircle\enspace #1}\\[0.7em]
        {\nohyf\Display\fontsize{19}{26}\selectfont\nh\color{ink} #2}};
  \end{tikzpicture}}
\begin{frame}[plain]\strut\end{frame}}}
